%%%%%%%%%%%%%%%%%%%%%%%%%%%%%%%%%%%%%%%%%
% The Legrand Orange Book
% LaTeX Template
% Version 2.4 (26/09/2018)
%
% This template was downloaded from:
% http://www.LaTeXTemplates.com
%
% Original author:
% Mathias Legrand (legrand.mathias@gmail.com) with modifications by:
% Vel (vel@latextemplates.com)
%
% License:
% CC BY-NC-SA 3.0 (http://creativecommons.org/licenses/by-nc-sa/3.0/)
%
% Compiling this template:
% This template uses biber for its bibliography and makeindex for its index.
% When you first open the template, compile it from the command line with the 
% commands below to make sure your LaTeX distribution is configured correctly:
%
% 1) pdflatex main
% 2) makeindex main.idx -s StyleInd.ist
% 3) biber main
% 4) pdflatex main x 2
%
% After this, when you wish to update the bibliography/index use the appropriate
% command above and make sure to compile with pdflatex several times 
% afterwards to propagate your changes to the document.
%
% This template also uses a number of packages which may need to be
% updated to the newest versions for the template to compile. It is strongly
% recommended you update your LaTeX distribution if you have any
% compilation errors.
%
% Important note:
% Chapter heading images should have a 2:1 width:height ratio,
% e.g. 920px width and 460px height.
%
%%%%%%%%%%%%%%%%%%%%%%%%%%%%%%%%%%%%%%%%%

%----------------------------------------------------------------------------------------
%	PACKAGES AND OTHER DOCUMENT CONFIGURATIONS
%----------------------------------------------------------------------------------------

\documentclass[11pt,fleqn]{book} % Default font size and left-justified equations

%%%%%%%%%%%%%%%%%%%%%%%%%%%%%%%%%%%%%%%%%
% The Legrand Orange Book
% Structural Definitions File
% Version 2.1 (26/09/2018)
%
% Original author:
% Mathias Legrand (legrand.mathias@gmail.com) with modifications by:
% Vel (vel@latextemplates.com)
% 
% This file was downloaded from:
% http://www.LaTeXTemplates.com
%
% License:
% CC BY-NC-SA 3.0 (http://creativecommons.org/licenses/by-nc-sa/3.0/)
%
%%%%%%%%%%%%%%%%%%%%%%%%%%%%%%%%%%%%%%%%%

%----------------------------------------------------------------------------------------
%	VARIOUS REQUIRED PACKAGES AND CONFIGURATIONS
%----------------------------------------------------------------------------------------

\usepackage{graphicx} % Required for including pictures
\graphicspath{{Pictures/}} % Specifies the directory where pictures are stored

\usepackage{lipsum} % Inserts dummy text

\usepackage{tikz} % Required for drawing custom shapes
\usetikzlibrary{shapes.geometric, arrows}

% tekið út \usepackage[english]{babel} % English language/hyphenation

\usepackage{enumitem} % Customize lists
\setlist{nolistsep} % Reduce spacing between bullet points and numbered lists

\usepackage{booktabs} % Required for nicer horizontal rules in tables

\usepackage{xcolor} % Required for specifying colors by name
\definecolor{ocre}{RGB}{34,139,34} % Define the orange color used for highlighting throughout the book

%----------------------------------------------------------------------------------------
%	MARGINS
%----------------------------------------------------------------------------------------

\usepackage{geometry} % Required for adjusting page dimensions and margins

\geometry{
	paper=a4paper, % Paper size, change to letterpaper for US letter size
	top=3cm, % Top margin
	bottom=3cm, % Bottom margin
	left=3cm, % Left margin
	right=3cm, % Right margin
	headheight=14pt, % Header height
	footskip=1.4cm, % Space from the bottom margin to the baseline of the footer
	headsep=10pt, % Space from the top margin to the baseline of the header
	%showframe, % Uncomment to show how the type block is set on the page
}

%----------------------------------------------------------------------------------------
%	FONTS
%----------------------------------------------------------------------------------------

\usepackage{avant} % Use the Avantgarde font for headings
%\usepackage{times} % Use the Times font for headings
\usepackage{mathptmx} % Use the Adobe Times Roman as the default text font together with math symbols from the Sym­bol, Chancery and Com­puter Modern fonts

\usepackage{microtype} % Slightly tweak font spacing for aesthetics
\usepackage[utf8]{inputenc} % Required for including letters with accents
\usepackage[T1]{fontenc} % Use 8-bit encoding that has 256 glyphs

%\usepackage[icelandic]{babel}

%----------------------------------------------------------------------------------------
%	BIBLIOGRAPHY AND INDEX
%----------------------------------------------------------------------------------------

%\usepackage[style=numeric,citestyle=numeric,sorting=nyt,sortcites=true,autopunct=true,babel=hyphen,hyperref=true,abbreviate=false,backref=true,backend=biber]{biblatex}
%\addbibresource{bibliography.bib} % BibTeX bibliography file
%\defbibheading{bibempty}{}

\usepackage{calc} % For simpler calculation - used for spacing the index letter headings correctly
\usepackage{makeidx} % Required to make an index
\makeindex % Tells LaTeX to create the files required for indexing

%----------------------------------------------------------------------------------------
%	MAIN TABLE OF CONTENTS
%----------------------------------------------------------------------------------------

\usepackage{titletoc} % Required for manipulating the table of contents

\contentsmargin{0cm} % Removes the default margin

% Part text styling (this is mostly taken care of in the PART HEADINGS section of this file)
\titlecontents{part}
	[0cm] % Left indentation
	{\addvspace{20pt}\bfseries} % Spacing and font options for parts
	{}
	{}
	{}

% Chapter text styling
\titlecontents{chapter}
	[1.25cm] % Left indentation
	{\addvspace{12pt}\large\sffamily\bfseries} % Spacing and font options for chapters
	{\color{ocre!60}\contentslabel[\Large\thecontentslabel]{1.25cm}\color{ocre}} % Formatting of numbered sections of this type
	{\color{ocre}} % Formatting of numberless sections of this type
	{\color{ocre!60}\normalsize\;\titlerule*[.5pc]{.}\;\thecontentspage} % Formatting of the filler to the right of the heading and the page number

% Section text styling
\titlecontents{section}
	[1.25cm] % Left indentation
	{\addvspace{3pt}\sffamily\bfseries} % Spacing and font options for sections
	{\contentslabel[\thecontentslabel]{1.25cm}} % Formatting of numbered sections of this type
	{} % Formatting of numberless sections of this type
	{\hfill\color{black}\thecontentspage} % Formatting of the filler to the right of the heading and the page number

% Subsection text styling
\titlecontents{subsection}
	[1.25cm] % Left indentation
	{\addvspace{1pt}\sffamily\small} % Spacing and font options for subsections
	{\contentslabel[\thecontentslabel]{1.25cm}} % Formatting of numbered sections of this type
	{} % Formatting of numberless sections of this type
	{\ \titlerule*[.5pc]{.}\;\thecontentspage} % Formatting of the filler to the right of the heading and the page number

% Figure text styling
\titlecontents{figure}
	[1.25cm] % Left indentation
	{\addvspace{1pt}\sffamily\small} % Spacing and font options for figures
	{\thecontentslabel\hspace*{1em}} % Formatting of numbered sections of this type
	{} % Formatting of numberless sections of this type
	{\ \titlerule*[.5pc]{.}\;\thecontentspage} % Formatting of the filler to the right of the heading and the page number

% Table text styling
\titlecontents{table}
	[1.25cm] % Left indentation
	{\addvspace{1pt}\sffamily\small} % Spacing and font options for tables
	{\thecontentslabel\hspace*{1em}} % Formatting of numbered sections of this type
	{} % Formatting of numberless sections of this type
	{\ \titlerule*[.5pc]{.}\;\thecontentspage} % Formatting of the filler to the right of the heading and the page number

%----------------------------------------------------------------------------------------
%	MINI TABLE OF CONTENTS IN PART HEADS
%----------------------------------------------------------------------------------------

% Chapter text styling
\titlecontents{lchapter}
	[0em] % Left indentation
	{\addvspace{15pt}\large\sffamily\bfseries} % Spacing and font options for chapters
	{\color{ocre}\contentslabel[\Large\thecontentslabel]{1.25cm}\color{ocre}} % Chapter number
	{}  
	{\color{ocre}\normalsize\sffamily\bfseries\;\titlerule*[.5pc]{.}\;\thecontentspage} % Page number

% Section text styling
\titlecontents{lsection}
	[0em] % Left indentation
	{\sffamily\small} % Spacing and font options for sections
	{\contentslabel[\thecontentslabel]{1.25cm}} % Section number
	{}
	{}

% Subsection text styling (note these aren't shown by default, display them by searchings this file for tocdepth and reading the commented text)
\titlecontents{lsubsection}
	[.5em] % Left indentation
	{\sffamily\footnotesize} % Spacing and font options for subsections
	{\contentslabel[\thecontentslabel]{1.25cm}}
	{}
	{}

%----------------------------------------------------------------------------------------
%	HEADERS AND FOOTERS
%----------------------------------------------------------------------------------------

\usepackage{fancyhdr} % Required for header and footer configuration

\pagestyle{fancy} % Enable the custom headers and footers

\renewcommand{\chaptermark}[1]{\markboth{\sffamily\normalsize\bfseries\chaptername\ \thechapter.\ #1}{}} % Styling for the current chapter in the header
\renewcommand{\sectionmark}[1]{\markright{\sffamily\normalsize\thesection\hspace{5pt}#1}{}} % Styling for the current section in the header

\fancyhf{} % Clear default headers and footers
\fancyhead[LE,RO]{\sffamily\normalsize\thepage} % Styling for the page number in the header
\fancyhead[LO]{\rightmark} % Print the nearest section name on the left side of odd pages
\fancyhead[RE]{\leftmark} % Print the current chapter name on the right side of even pages
%\fancyfoot[C]{\thepage} % Uncomment to include a footer

\renewcommand{\headrulewidth}{0.5pt} % Thickness of the rule under the header

\fancypagestyle{plain}{% Style for when a plain pagestyle is specified
	\fancyhead{}\renewcommand{\headrulewidth}{0pt}%
}

% Removes the header from odd empty pages at the end of chapters
\makeatletter
\renewcommand{\cleardoublepage}{
\clearpage\ifodd\c@page\else
\hbox{}
\vspace*{\fill}
\thispagestyle{empty}
\newpage
\fi}

%----------------------------------------------------------------------------------------
%	THEOREM STYLES
%----------------------------------------------------------------------------------------

\usepackage{amsmath,amsfonts,amssymb,amsthm} % For math equations, theorems, symbols, etc

\newcommand{\intoo}[2]{\mathopen{]}#1\,;#2\mathclose{[}}
\newcommand{\ud}{\mathop{\mathrm{{}d}}\mathopen{}}
\newcommand{\intff}[2]{\mathopen{[}#1\,;#2\mathclose{]}}
\renewcommand{\qedsymbol}{$\blacksquare$}
\newtheorem{notation}{Notation}[chapter]

% Boxed/framed environments
\newtheoremstyle{ocrenumbox}% Theorem style name
{0pt}% Space above
{0pt}% Space below
{\normalfont}% Body font
{}% Indent amount
{\small\bf\sffamily\color{ocre}}% Theorem head font
{\;}% Punctuation after theorem head
{0.25em}% Space after theorem head
{\small\sffamily\color{ocre}\thmname{#1}\nobreakspace\thmnumber{\@ifnotempty{#1}{}\@upn{#2}}% Theorem text (e.g. Theorem 2.1)
\thmnote{\nobreakspace\the\thm@notefont\sffamily\bfseries\color{black}---\nobreakspace#3.}} % Optional theorem note

\newtheoremstyle{blacknumex}% Theorem style name
{5pt}% Space above
{5pt}% Space below
{\normalfont}% Body font
{} % Indent amount
{\small\bf\sffamily}% Theorem head font
{\;}% Punctuation after theorem head
{0.25em}% Space after theorem head
{\small\sffamily{\tiny\ensuremath{\blacksquare}}\nobreakspace\thmname{#1}\nobreakspace\thmnumber{\@ifnotempty{#1}{}\@upn{#2}}% Theorem text (e.g. Theorem 2.1)
\thmnote{\nobreakspace\the\thm@notefont\sffamily\bfseries---\nobreakspace#3.}}% Optional theorem note

\newtheoremstyle{blacknumbox} % Theorem style name
{0pt}% Space above
{0pt}% Space below
{\normalfont}% Body font
{}% Indent amount
{\small\bf\sffamily}% Theorem head font
{\;}% Punctuation after theorem head
{0.25em}% Space after theorem head
{\small\sffamily\thmname{#1}\nobreakspace\thmnumber{\@ifnotempty{#1}{}\@upn{#2}}% Theorem text (e.g. Theorem 2.1)
\thmnote{\nobreakspace\the\thm@notefont\sffamily\bfseries---\nobreakspace#3.}}% Optional theorem note

% Non-boxed/non-framed environments
\newtheoremstyle{ocrenum}% Theorem style name
{5pt}% Space above
{5pt}% Space below
{\normalfont}% Body font
{}% Indent amount
{\small\bf\sffamily\color{ocre}}% Theorem head font
{\;}% Punctuation after theorem head
{0.25em}% Space after theorem head
{\small\sffamily\color{ocre}\thmname{#1}\nobreakspace\thmnumber{\@ifnotempty{#1}{}\@upn{#2}}% Theorem text (e.g. Theorem 2.1)
\thmnote{\nobreakspace\the\thm@notefont\sffamily\bfseries\color{black}---\nobreakspace#3.}} % Optional theorem note
\makeatother

% Defines the theorem text style for each type of theorem to one of the three styles above
\newcounter{dummy} 
\numberwithin{dummy}{section}
\theoremstyle{ocrenumbox}
\newtheorem{theoremeT}[dummy]{Theorem}
\newtheorem{itarefniutl}{Ítarefni}[chapter] %breytt
\newtheorem{problem}{Problem}[chapter]
\newtheorem{exerciseT}{Æfing}[chapter] %breytt
\theoremstyle{blacknumex}
\newtheorem{exampleT}{Example}[chapter]
\theoremstyle{blacknumbox}
\newtheorem{vocabulary}{Vocabulary}[chapter]
\newtheorem{definitionT}{Definition}[section]
\newtheorem{corollaryT}[dummy]{Corollary}
\theoremstyle{ocrenum}
\newtheorem{proposition}[dummy]{Proposition}

%----------------------------------------------------------------------------------------
%	DEFINITION OF COLORED BOXES
%----------------------------------------------------------------------------------------

\RequirePackage[framemethod=default]{mdframed} % Required for creating the theorem, definition, exercise and corollary boxes

% Theorem box
\newmdenv[skipabove=7pt,
skipbelow=7pt,
backgroundcolor=black!5,
linecolor=ocre,
innerleftmargin=5pt,
innerrightmargin=5pt,
innertopmargin=5pt,
leftmargin=0cm,
rightmargin=0cm,
innerbottommargin=5pt]{tBox}

% Exercise box	  
\newmdenv[skipabove=7pt,
skipbelow=7pt,
rightline=false,
leftline=true,
topline=false,
bottomline=false,
%backgroundcolor=ocre!10,
linecolor=ocre,
innerleftmargin=5pt,
innerrightmargin=5pt,
innertopmargin=5pt,
innerbottommargin=5pt,
leftmargin=0cm,
rightmargin=0cm,
linewidth=2pt]{eBox}	

% Definition box
\newmdenv[skipabove=7pt,
skipbelow=7pt,
rightline=false,
leftline=true,
topline=false,
bottomline=false,
linecolor=ocre,
innerleftmargin=5pt,
innerrightmargin=5pt,
innertopmargin=0pt,
leftmargin=0cm,
rightmargin=0cm,
linewidth=4pt,
innerbottommargin=0pt]{dBox}	

% Corollary box
\newmdenv[skipabove=7pt,
skipbelow=7pt,
rightline=false,
leftline=true,
topline=false,
bottomline=false,
linecolor=gray,
backgroundcolor=black!5,
innerleftmargin=5pt,
innerrightmargin=5pt,
innertopmargin=5pt,
leftmargin=0cm,
rightmargin=0cm,
linewidth=4pt,
innerbottommargin=5pt]{cBox}

% Creates an environment for each type of theorem and assigns it a theorem text style from the "Theorem Styles" section above and a colored box from above
\newenvironment{theorem}{\begin{tBox}\begin{theoremeT}}{\end{theoremeT}\end{tBox}}
\newenvironment{exercise}{\begin{eBox}\begin{exerciseT}}{\hfill{\color{ocre}\tiny\ensuremath{\blacksquare}}\end{exerciseT}\end{eBox}}				  
\newenvironment{definition}{\begin{dBox}\begin{definitionT}}{\end{definitionT}\end{dBox}}	
\newenvironment{example}{\begin{exampleT}}{\hfill{\tiny\ensuremath{\blacksquare}}\end{exampleT}}		
\newenvironment{corollary}{\begin{cBox}\begin{corollaryT}}{\end{corollaryT}\end{cBox}}	

%----------------------------------------------------------------------------------------
%	REMARK ENVIRONMENT
%----------------------------------------------------------------------------------------

\newenvironment{remark}{\par\vspace{10pt}\small % Vertical white space above the remark and smaller font size
\begin{list}{}{
\leftmargin=35pt % Indentation on the left
\rightmargin=25pt}\item\ignorespaces % Indentation on the right
\makebox[-2.5pt]{\begin{tikzpicture}[overlay]
\node[draw=ocre!60,line width=1pt,circle,fill=ocre!25,font=\sffamily\bfseries,inner sep=2pt,outer sep=0pt] at (-15pt,0pt){\textcolor{ocre}{R}};\end{tikzpicture}} % Orange R in a circle
\advance\baselineskip -1pt}{\end{list}\vskip5pt} % Tighter line spacing and white space after remark

%----------------------------------------------------------------------------------------
%	SECTION NUMBERING IN THE MARGIN
%----------------------------------------------------------------------------------------

\makeatletter
\renewcommand{\@seccntformat}[1]{\llap{\textcolor{ocre}{\csname the#1\endcsname}\hspace{1em}}}                    
\renewcommand{\section}{\@startsection{section}{1}{\z@}
{-4ex \@plus -1ex \@minus -.4ex}
{1ex \@plus.2ex }
{\normalfont\large\sffamily\bfseries}}
\renewcommand{\subsection}{\@startsection {subsection}{2}{\z@}
{-3ex \@plus -0.1ex \@minus -.4ex}
{0.5ex \@plus.2ex }
{\normalfont\sffamily\bfseries}}
\renewcommand{\subsubsection}{\@startsection {subsubsection}{3}{\z@}
{-2ex \@plus -0.1ex \@minus -.2ex}
{.2ex \@plus.2ex }
{\normalfont\small\sffamily\bfseries}}                        
\renewcommand\paragraph{\@startsection{paragraph}{4}{\z@}
{-2ex \@plus-.2ex \@minus .2ex}
{.1ex}
{\normalfont\small\sffamily\bfseries}}

%----------------------------------------------------------------------------------------
%	PART HEADINGS
%----------------------------------------------------------------------------------------

% Numbered part in the table of contents
\newcommand{\@mypartnumtocformat}[2]{%
	\setlength\fboxsep{0pt}%
	\noindent\colorbox{ocre!20}{\strut\parbox[c][.7cm]{\ecart}{\color{ocre!70}\Large\sffamily\bfseries\centering#1}}\hskip\esp\colorbox{ocre!40}{\strut\parbox[c][.7cm]{\linewidth-\ecart-\esp}{\Large\sffamily\centering#2}}%
}

% Unnumbered part in the table of contents
\newcommand{\@myparttocformat}[1]{%
	\setlength\fboxsep{0pt}%
	\noindent\colorbox{ocre!40}{\strut\parbox[c][.7cm]{\linewidth}{\Large\sffamily\centering#1}}%
}

\newlength\esp
\setlength\esp{4pt}
\newlength\ecart
\setlength\ecart{1.2cm-\esp}
\newcommand{\thepartimage}{}%
\newcommand{\partimage}[1]{\renewcommand{\thepartimage}{#1}}%
\def\@part[#1]#2{%
\ifnum \c@secnumdepth >-2\relax%
\refstepcounter{part}%
\addcontentsline{toc}{part}{\texorpdfstring{\protect\@mypartnumtocformat{\thepart}{#1}}{\partname~\thepart\ ---\ #1}}
\else%
\addcontentsline{toc}{part}{\texorpdfstring{\protect\@myparttocformat{#1}}{#1}}%
\fi%
\startcontents%
\markboth{}{}%
{\thispagestyle{empty}%
\begin{tikzpicture}[remember picture,overlay]%
\node at (current page.north west){\begin{tikzpicture}[remember picture,overlay]%	
\fill[ocre!20](0cm,0cm) rectangle (\paperwidth,-\paperheight);
\node[anchor=north] at (4cm,-3.25cm){\color{ocre!40}\fontsize{220}{100}\sffamily\bfseries\thepart}; 
\node[anchor=south east] at (\paperwidth-1cm,-\paperheight+1cm){\parbox[t][][t]{8.5cm}{
\printcontents{l}{0}{\setcounter{tocdepth}{1}}% The depth to which the Part mini table of contents displays headings; 0 for chapters only, 1 for chapters and sections and 2 for chapters, sections and subsections
}};
\node[anchor=north east] at (\paperwidth-1.5cm,-3.25cm){\parbox[t][][t]{15cm}{\strut\raggedleft\color{white}\fontsize{30}{30}\sffamily\bfseries#2}};
\end{tikzpicture}};
\end{tikzpicture}}%
\@endpart}
\def\@spart#1{%
\startcontents%
\phantomsection
{\thispagestyle{empty}%
\begin{tikzpicture}[remember picture,overlay]%
\node at (current page.north west){\begin{tikzpicture}[remember picture,overlay]%	
\fill[ocre!20](0cm,0cm) rectangle (\paperwidth,-\paperheight);
\node[anchor=north east] at (\paperwidth-1.5cm,-3.25cm){\parbox[t][][t]{15cm}{\strut\raggedleft\color{white}\fontsize{30}{30}\sffamily\bfseries#1}};
\end{tikzpicture}};
\end{tikzpicture}}
\addcontentsline{toc}{part}{\texorpdfstring{%
\setlength\fboxsep{0pt}%
\noindent\protect\colorbox{ocre!40}{\strut\protect\parbox[c][.7cm]{\linewidth}{\Large\sffamily\protect\centering #1\quad\mbox{}}}}{#1}}%
\@endpart}
\def\@endpart{\vfil\newpage
\if@twoside
\if@openright
\null
\thispagestyle{empty}%
\newpage
\fi
\fi
\if@tempswa
\twocolumn
\fi}

%----------------------------------------------------------------------------------------
%	CHAPTER HEADINGS
%----------------------------------------------------------------------------------------

% A switch to conditionally include a picture, implemented by Christian Hupfer
\newif\ifusechapterimage
\usechapterimagetrue
\newcommand{\thechapterimage}{}%
\newcommand{\chapterimage}[1]{\ifusechapterimage\renewcommand{\thechapterimage}{#1}\fi}%
\newcommand{\autodot}{.}
\def\@makechapterhead#1{%
{\parindent \z@ \raggedright \normalfont
\ifnum \c@secnumdepth >\m@ne
\if@mainmatter
\begin{tikzpicture}[remember picture,overlay]
\node at (current page.north west)
{\begin{tikzpicture}[remember picture,overlay]
\node[anchor=north west,inner sep=0pt] at (0,0) {\ifusechapterimage\includegraphics[width=\paperwidth]{\thechapterimage}\fi};
\draw[anchor=west] (\Gm@lmargin,-9cm) node [line width=2pt,rounded corners=15pt,draw=ocre,fill=white,fill opacity=0.5,inner sep=15pt]{\strut\makebox[22cm]{}};
\draw[anchor=west] (\Gm@lmargin+.3cm,-9cm) node {\huge\sffamily\bfseries\color{black}\thechapter\autodot~#1\strut};
\end{tikzpicture}};
\end{tikzpicture}
\else
\begin{tikzpicture}[remember picture,overlay]
\node at (current page.north west)
{\begin{tikzpicture}[remember picture,overlay]
\node[anchor=north west,inner sep=0pt] at (0,0) {\ifusechapterimage\includegraphics[width=\paperwidth]{\thechapterimage}\fi};
\draw[anchor=west] (\Gm@lmargin,-9cm) node [line width=2pt,rounded corners=15pt,draw=ocre,fill=white,fill opacity=0.5,inner sep=15pt]{\strut\makebox[22cm]{}};
\draw[anchor=west] (\Gm@lmargin+.3cm,-9cm) node {\huge\sffamily\bfseries\color{black}#1\strut};
\end{tikzpicture}};
\end{tikzpicture}
\fi\fi\par\vspace*{270\p@}}}

%-------------------------------------------

\def\@makeschapterhead#1{%
\begin{tikzpicture}[remember picture,overlay]
\node at (current page.north west)
{\begin{tikzpicture}[remember picture,overlay]
\node[anchor=north west,inner sep=0pt] at (0,0) {\ifusechapterimage\includegraphics[width=\paperwidth]{\thechapterimage}\fi};
\draw[anchor=west] (\Gm@lmargin,-9cm) node [line width=2pt,rounded corners=15pt,draw=ocre,fill=white,fill opacity=0.5,inner sep=15pt]{\strut\makebox[22cm]{}};
\draw[anchor=west] (\Gm@lmargin+.3cm,-9cm) node {\huge\sffamily\bfseries\color{black}#1\strut};
\end{tikzpicture}};
\end{tikzpicture}
\par\vspace*{270\p@}}
\makeatother

%----------------------------------------------------------------------------------------
%	LINKS
%----------------------------------------------------------------------------------------
\usepackage{float}
\usepackage{hyperref}
\hypersetup{hidelinks,backref=true,pagebackref=true,hyperindex=true,colorlinks=false,breaklinks=true,urlcolor=ocre,bookmarks=true,bookmarksopen=false}

\usepackage{bookmark}
\bookmarksetup{
open,
numbered,
addtohook={%
\ifnum\bookmarkget{level}=0 % chapter
\bookmarksetup{bold}%
\fi
\ifnum\bookmarkget{level}=-1 % part
\bookmarksetup{color=ocre,bold}%
\fi
}
}

%--------------------
% Viðbætur frá mér
%--------------------


\usepackage{todonotes}

%\usepackage[icelandic]{babel}

\usepackage{listings}

\definecolor{dkgreen}{RGB}{0,0.6,0}
\definecolor{gray}{RGB}{0.5,0.5,0.5}
\definecolor{mauve}{RGB}{0.58,0,0.82}
\definecolor{codegreen}{rgb}{0,0.6,0}
\definecolor{codegray}{rgb}{0.5,0.5,0.5}
\definecolor{codepurple}{rgb}{0.58,0,0.82}
\definecolor{backcolour}{rgb}{0.95,0.97,0.95}
\definecolor{arsenic}{rgb}{0.85, 0.90, 0.85}
\definecolor{brightgreen}{rgb}{0.4, 1.0, 0.0}

\lstdefinestyle{venjulegt}{
	language=Python,
	aboveskip=3mm,
	belowskip=1mm,
	showstringspaces=false,
	columns=flexible,
	basicstyle={\small\ttfamily},
	numbers=left,
	numberstyle=\tiny\color{gray},
	keywordstyle=\color{blue},
	backgroundcolor=\color{backcolour},   
	commentstyle=\color{red},
	keywordstyle=\color{codegreen},
	stringstyle=\color{codepurple},
	commentstyle=\color{dkgreen},
	stringstyle=\color{mauve},
	basicstyle=\fontfamily{cmtt}\footnotesize,
	keepspaces=true,
	breaklines=true,
	breakatwhitespace=true,
	tabsize=3,
	frame=tb,
	literate=%
	{Á}{{\'A}}1
	{á}{{\'a}}1
	{Ð}{{\DH}}1
	{ð}{{\dh}}1
	{É}{{\'E}}1
	{é}{{\'e}}1
	{í}{{\'i}}1
	{Í}{{\'I}}1
	{Ó}{{\'O}}1
	{ó}{{\'o}}1
	{Ú}{{\'U}}1
	{ú}{{\'u}}1
	{Ý}{{\'Y}}1
	{ý}{{\'y}}1
	{Þ}{{\TH}}1
	{þ}{{\th}}1
	{Æ}{{\AE}}1
	{æ}{{\ae}}1
	{Ö}{{\"O}}1
	{ö}{{\"o}}1}

\lstdefinestyle{uttak}{
	language=Python,
	aboveskip=1mm,
	belowskip=3mm,
	showstringspaces=false,
	columns=flexible,
	basicstyle={\small\ttfamily},
	numbers=left,
	numberstyle=\tiny\color{gray},
	backgroundcolor=\color{arsenic},
	keywordstyle=\color{gray},   
	commentstyle=\color{gray},
	keywordstyle=\color{gray},
	stringstyle=\color{gray},
	commentstyle=\color{gray},
	stringstyle=\color{gray}, 
	framexleftmargin=5mm, 
	frame=shadowbox, 
	rulecolor=\color{ocre},
	rulesepcolor=\color{ocre},
	basicstyle=\fontfamily{cmtt}\footnotesize,
	keepspaces=true,
	breaklines=true,
	breakatwhitespace=true,
	tabsize=3,
	frame=tb,
	literate=%
	{Á}{{\'A}}1
	{á}{{\'a}}1
	{Ð}{{\DH}}1
	{ð}{{\dh}}1
	{É}{{\'E}}1
	{é}{{\'e}}1
	{í}{{\'i}}1
	{Í}{{\'I}}1
	{Ó}{{\'O}}1
	{ó}{{\'o}}1
	{Ú}{{\'U}}1
	{ú}{{\'u}}1
	{Ý}{{\'Y}}1
	{ý}{{\'y}}1
	{Þ}{{\TH}}1
	{þ}{{\th}}1
	{Æ}{{\AE}}1
	{æ}{{\ae}}1
	{Ö}{{\"O}}1
	{ö}{{\"o}}1}

\lstset{style=venjulegt}

\renewcommand{\lstlistingname}{Kóðabútur}

\newcommand{\comment}[1]{}

\usepackage[lastexercise,answerdelayed]{exercise}
\newboolean{firstanswerofthechapter}  
\renewcommand{\AnswerName}{Æfing}
\renewcommand{\AnswerHeader}{\ifthenelse{\boolean{firstanswerofthechapter}}%
	{\bigskip\noindent\textcolor{ocre}{\textbf{\huge Kafli \thechapter}}\newline\newline%
		\noindent\bfseries{\textcolor{ocre}{\AnswerName\ \ExerciseHeaderNB}}\smallskip\newline}
	{\noindent\bfseries{\textcolor{ocre}{\AnswerName\ \ExerciseHeaderNB}}\smallskip\newline}}
\setlength{\QuestionIndent}{16pt}

% Exercise box	  
\newmdenv[skipabove=7pt,
skipbelow=7pt,
outerlinewidth=0,
rightline=true,
leftline=true,
topline=true,
bottomline=true,
backgroundcolor=ocre!10,
linecolor=ocre!65,
innerleftmargin=15pt,
innerrightmargin=15pt,
innertopmargin=10pt,
innerbottommargin=10pt,
leftmargin=0cm,
rightmargin=0cm,
roundcorner=50,
fontcolor=gray!85,
linewidth=4pt]{valBox}

\newcounter{valBox}[section]
\newenvironment{itarefni}{\begin{valBox}\begin{itarefniutl}}{\end{itarefniutl}\end{valBox}}

\usepackage{tasks} % Insert the commands.tex file which contains the majority of the structure behind the template

%\hypersetup{pdftitle={Title},pdfauthor={Author}} % Uncomment and fill out to include PDF metadata for the author and title of the book

%----------------------------------------------------------------------------------------

\begin{document}

%----------------------------------------------------------------------------------------
%	TITLE PAGE
%----------------------------------------------------------------------------------------

\begingroup
\thispagestyle{empty} % Suppress headers and footers on the title page
\begin{tikzpicture}[remember picture,overlay]
\node[inner sep=0pt] (background) at (current page.center) {\includegraphics[width=\paperwidth]{background.pdf}};
\draw (current page.center) node [fill=ocre!30!white,fill opacity=0.6,text opacity=1,inner sep=1cm]{\Huge\centering\bfseries\sffamily\parbox[c][][t]{\paperwidth}{\centering Inngangur að forritun í Python\\[15pt] % Book title
{\Large Forritun fyrir byrjendur}\\[20pt] % Subtitle
{\huge Valborg Sturludóttir}}}; % Author name
\end{tikzpicture}
\vfill
\endgroup

%----------------------------------------------------------------------------------------
%	COPYRIGHT PAGE
%----------------------------------------------------------------------------------------

\newpage
~\vfill
\thispagestyle{empty}

\noindent Copyright \copyright\ 2020 Valborg Sturludóttir\\ % Copyright notice

% \noindent \textsc{Published by Publisher}\\ % Publisher

% \noindent \textsc{book-website.com}\\ % URL

\noindent Licensed under the Creative Commons Attribution-NonCommercial 3.0 Unported License (the ``License''). You may not use this file except in compliance with the License. You may obtain a copy of the License at \url{http://creativecommons.org/licenses/by-nc/3.0}. Unless required by applicable law or agreed to in writing, software distributed under the License is distributed on an \textsc{``as is'' basis, without warranties or conditions of any kind}, either express or implied. See the License for the specific language governing permissions and limitations under the License.\\ % License information, replace this with your own license (if any)

\noindent \textit{First printing, Sept 2020} % Printing/edition date

%----------------------------------------------------------------------------------------
%	TABLE OF CONTENTS
%----------------------------------------------------------------------------------------

%\usechapterimagefalse % If you don't want to include a chapter image, use this to toggle images off - it can be enabled later with \usechapterimagetrue

\chapterimage{chapter_head_1.pdf} % Table of contents heading image

\pagestyle{empty} % Disable headers and footers for the following pages

\tableofcontents % Print the table of contents itself

\cleardoublepage % Forces the first chapter to start on an odd page so it's on the right side of the book

\pagestyle{fancy} % Enable headers and footers again

%----------------------------------------------------------------------------------------
%	PART
%----------------------------------------------------------------------------------------

\part{Fyrsti hluti}

\chapterimage{chapters0.png} % Chapter heading image



\chapter{Inngangur}

\section{Tilgangur bókarinnar}\index{Tilgangur Python}

Þessi bók fjallar um þau undirstöðu atriði sem þarf að kynna til að ná tökum á forritun í Python. 
Höfundi finnst mikilvægt að kenna námsefnið með íslenskum hugtökum þar sem ætlunin er að nota hana í kennslu í íslenskum framhaldsskólum. 
Ef nemendur ætla að leggja fyrir sig tölvunarfræði í framhaldsnámi er nauðsynlegt að búa yfir ríkulegu íðorðasafni, þess þá heldur ef nemandi hyggst framfleyta fræðunum. 
Hugtök verða þó líka sett fram á ensku því lesandi gæti óskað að fletta upp ítarefni sem meira er til af á netinu á ensku en íslensku.

Það er algengur misskilningur að forritarar kunni rosalega mörg forritunarmál, eins og fólk sem getur talað mörg tungumál, eða það að kunna rosalega mörg mál geri þig að góðum forritara.
Þvert á móti.
Að sýna hæfni og leikni í einu máli er auðveldlega yfirfæranlegt á önnur mál sé þess þörf.
Þess vegna er út í hött að spyrja: ,,hvað kanntu mörg forritunarmál?''
Ekki aðeins eru tungumál og forritunarmál gerólík, forritunarmál eru formleg mál og þekking á einu hlutbundnu máli er nær því að vera jafn frábrugðið öðru í grunninn eins og mállýskur innan tungumála.
Nær væri að spyrja hvort viðkomandi hafi meiri áhuga á framenda eða bakenda forritun, hvert er skemmtilegasta reikniritið sem viðkomandi hefur útfært eða hvert er það forritunarmál sem viðkomandi grípur oftast í.

Einnig er það algengur misskilningur að það fyrsta sem fólk gerir er að búa til tölvuleik.
Það þarf mikla undirstöðu kunnáttu til þess að geta búið til tölvuleiki, alveg eins og áður en hafist er handa við að skrifa bók þarf að læra stafrófið.
Þessi grunnvinna finnst mörgum vera leiðigjörn.
Að mati höfundar er það vegna þess að við erum svo vön því að nota tölvur dagsdaglega, svo fræðigreinin sem tæknin byggir á hlýtur líka að vera okkur kunnug ekki satt?
Nei, alveg eins og dýralækningar eru okkur ekki augljósar við að eiga gæludýr og pípulagnir heldur ekki við að nota klósett.
Innan tölvunnar eru ákveðnar grunneiningar sem eru notandanum ekki augljósar, af góðri ástæðu, það væri hrikalegt ef við þyrftum öll að vera píparar til þess að geta notað klósett.
Þó að þessi samlíking hafi verið heldur gróf þá sýnir hún að það eru svo margir hlutar sem eru okkur huldir að við hreinlega vitum ekki hvað við vitum ekki.
Því er nauðsynlegt að læra grunninn vel og fara rólega yfir hann svo þegar við ætlum að fara að afrita og líma kóða, frá síðum eins og stackoverflow, þá vitum við allavega hvað sá kóði gerir (nokkurn veginn).

Það er gífurlega mikilvægt að kunna að nota tólin sem við nýtum okkur til þess að við getum nýtt þau skynsamlega.
Fartölvur í dag eru mun öflugri en tölvan sem kom fólki á tunglið upphaflega en þó ósennilegt að margar þeirra séu nýttar í eins flókna hluti og geimferðir.
Hins vegar lendum við oft í því að þurfa að leysa einhver tiltekin verkefni og ef við gætum nýtt okkur vélaraflið betur værum við sneggri að því.

Það síðasta sem er þess virði að taka fram er mikilvægi opins (e. open source) og frjáls (e. free) hugbúnaðar.
Þessi bók er gefin út undir Creative Commons 3.0 sem þýðir að hver sem er má afrita hana, prenta hana og nota með því skilyrði að afleidd verk vísi til höfundar en að óheimils sé að hagnast efninu.

Internetið eins og við þekkjum það myndi ekki keyra ef það væri ekki fyrir framlög forritara sem viðhalda opnum hugbúnaði.
Stefnan í heiminum í dag virðist þó vera í áttina að lokuðum réttindavörðum hugbúnaði sem er mjög miður. 

Það að hugbúnaður sé opinn þýði að hver sem er geti skoðað kóðann á bak við hann, gert breytingar og deilt með öðrum.
Hugmyndin á bak við frjálsan hugbúnað er svipuð nema lagt er upp með þá hugmynd að það sé siðferðisleg skylda fólks að hugbúnaður sé sem aðgengilegastur því frelsi í hugbúnaði skilar sér í frjálsara samfélagi.

\section{Hvers vegna Python?}\index{Hvers vegna Python?}

Ástæður þess að Python er gott mál til þess að byrja á að skoða eru eftirfarandi: \footnote{Strax í þessum texta koma fyrir hugtök sem verða skýrð betur seinna, ekki missa kjarkinn.}.

\begin{enumerate}
	\item \textit{Málskipanin} er mjög svipuð mannlegu máli svo það ætti að vera auðvelt að læra hvernig eigi að ,,tala'' við tölvuna.
	\item  Python er \textit{kvikt tagað} forritunarmál, það þýðir að notandinn þarf ekki að gefa upp hvers konar \textit{gagnatýpur} er unnið með. 
	Þetta gerir það að verkum að notandinn þarf ekki að læra urmull af lykilorðum áður en byrjað er að forrita.
	\item  Python er ekki alveg \textit{hlutbundið} forritunarmál, sem gerir það að verkum að notandinn þarf ekki að læra hvernig á að beita hlutbundinni forritun fyrr en góð undirstaða er þegar komin.
	\item Python er frítt og aðgengilegt öllum helstu stýrikerfum og einnig er hægt að forrita yfir netið í vafra og því óþarfi fyrir notandann að setja nokkuð upp sé þess óskað.
	\item Python er æðra forritunarmál (e. high level programming language), ekki það að margir séu að kenna assembly í dag.
	\item Python er mikið notað, algengt mál svo það er praktískt að hafa undirstöðu skilning á því.
	\item Nefnt í höfuð á Monty Pyhton grínhópsins.
\end{enumerate}

\section{Uppsetning}\index{Uppsetning}
Uppbygging bókarinnar er þannig að fyrri hlutinn snýr að því að kynna lesandann fyrir grunnvirkni Python; málskipan, lykilhugtök og lykilorð, gagnatýpur, lykkjur og föll.
Seinni hlutinn snýr svo að því að beita þekkingu úr fyrri hlutanum í hlutbundinni forritun, þar verða kynntir til sögunnar klasar og aðferðir sem lesandinn útfærir upp á eigin spýtur.
Í lok hvers kafla eru svo verkefni til að reyna á leikni lesandans, lausnir við þeim öllum má finna í lausnarhluta aftast í bókinni.

Ekki er búist við neinni fyrri kunnáttu við lestur þessarar bókar, hún á að geta staðið fyrir sínu án þess að lesandinn búi yfir nokkurri þekkingu á sviði tölvunarfræði eða forritunar.
Ef slík þekking er fyrir hendi gæti lesandanum þótt ágætt að fara hratt í gegnum fyrri hluta bókarinnar og einbeita sér að verkefnum úr seinni hlutanum.

Víðs vegar um bókina má finna númeraða kóðabúta.
Ástæðan fyrir því er að þessum kóðabútum er auðveldara að viðhalda heldur en skjáskotum úr vinnubókum og því er heldur vísað í bækur sem eru aðgengilegar lesendum og frumstæðari framsetning ræður heldur ríkjum hér.
\lstset{style=uttak}
\lstset{style=venjulegt}
\begin{lstlisting}[caption=Kóðabútar kynntir til sögunnar]
# Svona líta kóðabútar út
# kóðann má allan afrita og keyra til að sjá þá virkni sem verið er að kynna
\end{lstlisting}
\lstset{style=uttak}
\begin{lstlisting}
# Svona lítur svarið út þegar kóðabúturinn fyrir ofan er keyrður
# Notið úttakið til að bera við eigin niðurstöður
\end{lstlisting}
\lstset{style=venjulegt}

Einnig eru á nokrum stöðum stuttar efnisgreinar af ítarefni sem er ekki nauðsynlegt að hafa fullan skilning á en þó gott að skoða, sérlega fyrir þá lesendur sem vilja leggja frekara nám í tölvunarfræðum fyrir sig.
Þær líta svona út:

\begin{itarefni}
\textbf{Titill á ítarefni}
Þennan texta má leiða hjá sér við flýtilestur en gott að hafa skilning á ef lesandi vill ná góðum tökum á efninu.
\end{itarefni}

\section{Að keyra kóða}\index{Að keyra kóða}\label{uk:keyra-koda}

Það fyrsta sem nemendur vilja yfirleitt gera er að byrja að skrifa sinn eigin kóða. 
Áður en við komumst svo langt þarf að útskýra hvernig það er gert. 
Þessi kennslubók byggir á notkun Jupyter Notebooks með hjálp Anaconda hugbúnaðarins, sem er öflugt pakkakerfi og tólakista sem hefur upp á mikið meira en bara Jupyter að bjóða. 
Hægt er að nálgast Anaconda á \href{www.anaconda.com}{anaconda.com}.
Hægt er að nota Jupyter án þess að ná í Anaconda með síðum eins og \href{www.cocalc.com}{cocalc.com}. 
Einnig er hægt að keyra kóða á netinu í gegnum síður eins og \href{www.repl.it}{repl.it}, nota ritla (eins og notepad eða sublime) og keyra .py skrár í skipanalínu, eða nota þyngri umhverfi eins og pycharm sem eru sérhönnuð fyrir hugbúnaðarþróun.

Jupyter er byrjendavænt umhverfi og því eru verkefnin sem fylgja þessari bók sett upp í Jupyter vinnubókum,sem lesandinn getur nýtt sér. 
Virkninni er skipt upp í sellur og keyrsluröð sellanna skiptir máli, við sjáum seinna mikilvægi þess að geta skipt upp kóða svona og hvers vegna þetta umhverfi er þægilegt til að byrja í. 
En hver sella hefur aðgang að svokölluðu skilgreiningarsvæði vinnubókarinnar en er þó sín eigin eining, því má keyra eina sellu í einu án þess að keyra allan kóðann í vinnubókinni.
Þetta er þægilegt ef við lendum í því að fá villu í einni sellu þá hefur ekki áhrif á neina aðra sellu og við getum haldið áfram.

En við keyrslu á kóða þarf einnig að hafa í huga að tölvan gerir nákvæmlega það sem við segjum henni að gera og ekkert annað.
Og þá komum við niður á stórt vandamál, að tölvur eru mjög bókstaflegar og vitlausar.
Þær skortir allt vit, þær reyna ekki að hafa vit fyrir okkur. 
Þær gera nákvæmlega það sem við biðjum um.
Nákvæmlega eins og við biðjum um það.

Þannig að ef ég ætlaði að segja tölvu að smyrja handa mér hnetusmjörs og sultu samloku þá þyrfti ég að segja vélinni að gera eftirfarandi í nákvæmlega þessari röð:
\vspace{0.4cm}
\begin{enumerate}
	\item taka fram hníf
	\item taka fram tvær brauðsneiðar
	\item opna hnetusmjörið
	\item setja beitta endann ofan í hnetusmjörið þannig að hann nái upp 50gr af hnetusmjöri
	\item setja hnetusmjörið sem er á hnífnum á miðja brauðsneiðina
	\item nota hnífinn til þess að smyrja hnetusmjörinu á þá hlið sem hnetusmjörið er nú þegar á, og enga aðra
	\item taka fram skeið
	\item opna sultuna
	\item setja kúpta enda skeiðarinnar ofan í sultukrukkuna 
	\item taka skeiðina upp úr sultukrukkunni með kúfaða skeið af sultu
	\item setja sultuna á hina brauðsneiðina
	\item nota skeiðina til að smyrja sultunni yfir þá hlið brauðsneiðarinnar sem sultan er á og enga aðra hlið
	\item  setja brauðsneiðarnar saman þannig að hnetusmjörið og sultan snertist og hornin mætast öll. 
\end{enumerate} 
\vspace{0.4cm}
Takið eftir að hér er gert ráð fyrir þó nokkru og ef vélin kann ekki nú þegar skil á eftirfarandi, mun þetta klúðrast: 
\vspace{0.2cm}
\begin{enumerate}
	\item taka fram
	\item hnífur
	\item opna
	\item mæla 50 gr
	\item smyrja
	\item hlið á brauðsneið
	\item miðja á brauðsneið
	\item skeið
	\item kúfað
\end{enumerate} 
\vspace{0.2cm}

Þessi útskýring á samlokugerð kann að vera alveg ofboðslega óþarflega nákvæm þá er ekki víst að úr þessu verði nokkur samloka.
Þetta könnumst við öll við, að tölvur gera það sem þeim er sagt, ekki það sem við viljum.

Helsta verkefni forritara er að búta niður verkefni í svo litla hluta að hægt er að útskýra þá fyrir tölvu.
Ekki búast við því að setjast niður við fyrsta verkefni og ætlast svo til að búa til tölvuleik eða hakka banka.
Forritun er einnig frábrugðin þeirri venjulegu tölvunotkun sem við höfum vanist dagsdaglega.
Dagsdaglega erum við ekki að gefa tölvunni okkar eigin skipanir heldur að beita skipunum sem aðrir forritara hafa samið og sett upp í hugbúnaðinn sem við erum að nota.

%------------------------------------------------

\section{Málskipan}\index{Málskipan}

\textbf{Málskipan} (e. syntax) er safn þeirra reglna sem við þurftum til þess að skrifa kóða sem tölvan skilur, svo að hann þýðist í vélamál, þær reglur sem við þurfum að fara eftir þegar við forritum, þær reglur sem forritunarmálið býst við að við förum eftir. 
Ef við brjótum þessar reglur fáum við villu, og einhver algengasta villa sem hægt er að fá er málskipanarvilla (e. syntax error). 
Python er frábrugðið öðrum forritunarmálum á þann hátt að málskipanin krefst þess að kóðinn sé settur upp á ákveðinn hátt. 
Líkja má því við að þurfa ekki að hafa greinarmerki í huga þegar við ljúkum setningum heldur setjum við orðin okkar á réttan stað í samræðum.

\subsection{Uppsetning á kóða}\index{Uppsetning á kóða}
Þessi kóðabútur er þannig uppsettur að allar línur byrja jafnlangt til vinstri, eins og hver setning í töluðu máli stendur hver lína fyrir sínu, ein og sér.
\begin{lstlisting}[caption=Réttur Python kóði, label=lst:inng-kóðadæmi]
# Réttur Python kóði sem keyrist
4 + 8
5 + 6
breyta = 9 * 2
\end{lstlisting}

Við munum sjá seinna hverjar reglurnar eru en sjáum dæmi sem, ef keyrt, veldur villu (á eftir villu) tekið er fram með athugasemdum hverjar villurnar verða.

Ekki hafa áhyggjur af því að sjá ekki endilega hvað það er sem er að, hér er einungis verið að sýna dæmi um slæma uppsetningu og illa skrifaðan kóða.

\begin{lstlisting}[caption=Rangur Python kóði sem veldur villu, label=lst:inng-malskipanarvilla]
# Illa skrifaður Python kóði sem keyrist ekki
4 + 8. # punkturinn veldur málskipanarvillu
	5 + 6 # inndrátturinn hér er rangur
breyta = 9 * 2 * # málskipanarvilla fæst hér því síðasta táknið er á röngum stað
\end{lstlisting}

Svona inndrætti er einungis beitt ef lína á beinlínis að hanga undir línunni að ofan og tilheyra henni, eins og kommusetning sem getur ekki staðið ein og sér og er háð samhengi setningarinnar sem hún tilheyrir.
Þess vegna þarf að huga að því hvernig Python kóði er uppsettur. 
Í öðrum málum eru notuð greinamerki til að segja tölvunni að lína sé búin og að aðrar línur eigi að heyra undir eitthvað ákveðið samhengi en ekki í Python, þar er treyst á að forritarinn setji kóðann upp á máta sem hægt er að sjá að sé réttur. 
Dæmi um hvernig línur geta verið aðgreindar í öðrum málum:

\begin{lstlisting}[language=Java , caption=Dæmi um annað mál sem er strangt tagað og með greinamerkjum]
// Dæmi um kóða í forritunarmálinu Java
int i = 7;
i + 5;

// Þetta myndi líka ganga í Java en ekki í Python
int i = 7; i + 5;
\end{lstlisting}

\begin{lstlisting}[language=Lisp, caption=Dæmi um annað mál sem byggir á afmörkuðu samhengi en með greinamerkjum]
; Lisp
(setq x 10)
(setq y 34.567)

(print x)
(print y)
\end{lstlisting}

Í þessum tveimur frábrugðnu málum sem voru tekin sem dæmi var óþarfi að setja kóðann í mismunandi línur, því greinamerkin væru nóg til að aðgreina hverja línu fyrir sig. 
Hins vegar er það góð venja að skrifa kóða sem er læsilegur öðru fólki. Í Java eru greinarmerkin semikommur (;) en í Lisp eru línur og samhengi afmörkuð með svigum. 
Python byggist hinsvegar á því að forritarinn stilli öllu upp rétt með réttum inndrætti. 

\subsection{Gagnatýpur og lykilorð}\index{Gagnatýpur og lykilorð}

Í Python eru nokkrar grunn gagnatýpur sem við munum kynnast í þessari bók. 
Ástæðan fyrir því að þær eru kallað grunntýpur er sú að þær fylgja með Python uppsetningunni og notandinn getur beitt þeim í samræmi við það sem þær eru færar um, sem má skoða í skjölun Python \href{https://www.python.org/doc/}{https://www.python.org/doc/}. 
Týpa eða tag er hugtak sem þýðir að hlutur sé af einhverri ákveðinni tegund sem má framkvæma ákveðnar aðgerðir á, í þessari bók verða týpur ýmist kallarðar það eða tög. 
Lesandi þekkir muninn á orðum og tölum úr daglegu tali og veit að hægt er að framkvæma mismunandi aðgerðir á þessum mismunandi týpum, eins og hægt er að skipta út hástöfum fyrir lágstafi í orðum en ekki tölum og hægt er að hefja tölur í veldi en ekki orð. 
Að sama skapi eru til aðgreinanlegar týpur sem tölvan kann skil á og leyfir ákveðnar aðgerðir á.
Í fyrri hluta þessarar bókar verða gerð skil á tveimur talnatýpum (heiltölum og fleytitölum), strengjum, listum, sanngildum, orðabókum (einnig kallaðar hakkatöflur), nd-um og mengjum.

Lykilorð eru orð sem eru frátekin og birtast þau græn í Jupyter vinnubók. 
Hver gagnatýpa hefur eitt lykilorð og eru einnig nokkur innbyggð föll í Python, sem við kynnumst fljótlega, með frátekin orð. 
Forðast skal að yfirskrifa þessi lykilorð, en gerist það þá er auðvelt að laga það í Jupyter. 
Hver vinnubók hefur sinn kjarna til að vinna á og það eina sem þarf að gera í aðstæðum þar sem innbyggt orð er allt í einu farið að þýða eitthvað annað þá dugir að endurræsa kjarnann.
Kjarninn í vinnubókinni er hvaða túlk eða þýðanda er verið að nota til þess að láta tölvuna skilja kóðann.
Í okkar tilfelli erum við að nota Python 3.


\chapter{Tölur og breytur}\index{Tölur og breytur}\label{k:tolur}
Í þessum kafla ætlum við að hefjast handa við að forrita. 
Það fyrsta sem við ætlum að gera er að kynnast talnatýpum og keyra kóða eins og við værum að nota reiknivél. 
Við könnumst við reiknivélar og hvernig þær afgreiða röð aðgerða. 
Nú viljum við sannreyna að þær reikniaðgerðir sem við þekkjum séu til í Python og að þegar við keyrum kóðann okkar þá verði útkoman sú sama og við áttum von á. 
Við viljum líka geta geymt útkomuna okkar til að nota aftur seinna, til þess þurfum við breytur (e. variables).

\section{Tölur - talnatýpur}\index{Tölur - talnatýpur}
Í Pyhton eru í grunninn tvær týpur af tölum (en til eru tvær týpur af hvorri fyrir sig, sem snýr meira að minnisnotkun og er út fyrir svið þessarar bókar). 
Þær eru:

\begin{itemize}
	\item \textbf{Heiltölur} - tölur sem eru ekki með neinum aukastaf. 
	Á ensku eru þessar tölur kallaðar integers og er lykilorð þeirra því \textbf{int}.
	\item \textbf{Fleytitölur} - tölur sem eru með aukastaf, sem er fyrir aftan punkt (ekki kommu, fleytitölur eru oft kallaðar kommutölur á íslensku). 
	Á ensku eru þessar tölur kallaðar floating point numbers og er því lykilorðið þeirra \textbf{float}.
\end{itemize}

\begin{lstlisting}[caption=Heiltölur og fleytitölur]
# Heiltölur, enginn aukastafur
42
100000
-139

# Fleytitölur, aukastafur/ir fyrir aftan punkt
4.0
3.1415926
-100.98
\end{lstlisting}

\section{Reikniaðgerðir og tákn}\index{Reikniaðgerðir og tákn}
Grunnreikniaðgerðir eru nokkrar sem við könnumst við úr grunnskóla en aðrar eru framandi og við skulum skoða aðeins betur.

Táknin eru flest eins og á reiknivélum +, -, *, / en þar að auki er annars konar deiling sem er táknuð með tveimur deilimerkjum //, veldishafning er táknuð með tveimur margföldunarmerkjum **, og svo er leifareikningur táknaður með \%.
Heiltöludeiling og leifarreikningur eru líklega ný á nálinni fyrir flestum lesendum og því allt í lagi að útskýra þær aðgerðir aðeins nánar.
Þessar aðgerðir eru einmitt mjög skyldar í raun.
Deilingin segir okkur hversu oft ein tala gengur upp í aðra þar sem útkoman er heil tala (eða fleytitala með 0 sem eina aukastafinn), okkur er sama um afganginn sem verður eftir.
Í þessari deilingu er svarið 2 við bæði 5//2 og 4//2.
En í leifareikningnum viljum við eingöngu vita hver er afgangurinn þegar heiltöludeilingu er beitt svo 5\%2 væri 1, því það er einn í afgang þegar fimm er deilt með tveimur.
Og það er 0 í afgang þegar fjórum er deilt með tveimur svo 4\%2 er 0.

Í eftirfarandi dæmum í kóðabút \ref{lst:reiknadg} er vert að draga fram nokkur atriði sem eru ekki augljós byrjanda. 
Það fyrsta er að myllumerkið (\#) þýðir að allt sem kemur fyrir aftan það er \textit{athugasemd}, athugasemdir eru engöngu til að gera kóða læsilegri fyrir fólk, þær eru hunsaðar af tölvunni þegar hún breytir kóðanum í eitthvað sem hún skilur.
Eins og sést í línu merktri númer 20 er athugasemdin svo löng að hún birtist okkur sem tvær línur en hún er í keyrslu tölvunnar álitin ein heild línu 20.
Þess vegna þurfum við ekki að hafa áhyggjur af þessum inndrætti sem birtist, hann er í rauninni ekki til staðar þar sem þessi hluti textans er ein heild.
Einnig eru þarna bil á milli talna fremst í línu og tákna, það er líka til að gera kóðan læsilegri, bilin mega bara ekki vera fremst í línunni enn sem komið er.
Athugasemdir í kóða eru mjög mikilvægur hluti af skjölun kóða og ættu öll sem vilja tileinka sér forritun að venja sig á að skrifa athugasemdir.
Í fyrstu erum við ekki að skrifa flókinn kóða svo athugasemdirnar segja okkur ekki mikið, en þegar kóðinn er ekki augljós eða lausn á verkefni ekki augljós er gott að skrifa athugasemdir.
Flest allir kóðabútar eru skjalaðir með athugasemdum til að gera þá læsilegri því allur kóði í bókinni er skrifaður fyrir fólk til að skilja.
Kóði sem þið komið til með að skrifa seinna meir á einnig að vera ykkur sjálfum skiljanlegur þegar þið komið að honum seinna.
Því er gott að venja sig strax á að skrifa lýsandi athugasemdir. 

\begin{lstlisting}[caption=Reikniaðgerðir, label=lst:reiknadg]
# Samlagning framkvæmd með + 
# Þegar eftirfarandi kóði er keyrður ætti útkonan að vera 10
6 + 4 

# Frádráttur framkvæmdur með -
# Þegar eftirfarandi kóði er keyrður ætti útkonan að vera 10
14 - 4 

# Margföldun framkvæmd með * 
# Þegar eftirfarandi kóði er keyrður ætti útkonan að vera 10
10 * 2 

# Deiling framkvæmd með / 
# Athugið að þetta er fleytitöludeiling sem skilar nákvæmu svari
# Þegar eftirfarandi kóði er keyrður ætti útkonan að vera 10.0
60 / 6 

# Heiltöludeiling framkvæmd með //
# Athugið að þessi deiling er frábrugðin þeirri sem þið kannist við
# Hér viljum við vita hversu oft, heil tala, ein tala gengur upp í aðra og okkur er sama um afganginn
# Þegar eftirfarandi kóði er keyrður ætti svarið að vera 10
177 // 17

# Veldishafning framkvæmd með **
# Hér er mikilvægt, eins og með deilinguna, að hafa í huga hvor talan kemur á undan.
# Fyrst kemur talan sem hefja á í veldi og svo kemur talan sem er veldisvísirinn
# Þegar eftirfarandi kóði er keyrður ætti svarið að vera 9
3 ** 2

# Leifareikningur framkvæmdur með % (e. modulus)
# Þetta er eitthvað alveg nýtt og framandi, en þó ekki óskiljanlegt
# Það sem þetta reiknar er hversu mikil leif eða afgangur er eftir þegar heiltöludeilingu er beitt.
# Þegar eftirfarandi kóði er keyrður ætti svarið að vera 7
177 % 17

\end{lstlisting}

\begin{valBox}
	Í öllum þessum dæmum var verið að vinna með heiltölur, þó var útkoman úr deilingunni (stundum kölluð fullkomin deiling) fleytitala. 
	Hvað gerist ef þessir sömu útreikningar eru gerðir með fleytitölum? 
	Ef við myndum skipta út hverri tölu fyrir sig og setja í staðinn sömu tölu með .0 fyrir aftan þá yrðu útkomurnar þær sömu nema fleytitölur. 
	En hvað gerist ef við breytum aðeins fyrri tölunni en ekki seinni tölunni?
	Þá ertu að nota ólíkar týpur og slíkt er vandmeðfarið, en í þessu tilviki er það í lagi þar sem Python gerir þá ráð fyrir að það sé í lagi að reikna allt með fleytitölum og framkvæmir reikninginn eins og þú hafir verið að beita fleytitölum í hvívetna og niðurstaðan verður þá að sjálfsögðu fleytitala.
\end{valBox}

\section{Breytur}\index{Breytur}
Nú höfum við séð hvernig má keyra kóða einfaldlega eins og í reiknivél.
Höldum okkur við samlíkinguna um reiknivélina til að útskýra breytur.
Á hefbundinni reiknivél sem notuð er í stærðfræðitíma í framhaldsskóla er takki sem á stendur ANS.
Það stendur fyrir answer og ef ýtt er á hann getur vélin geymt síðasta gildið sem hún gaf sem svar og unnið svo með það til að gefa næsta svar.
Flottari vélar geta svo geymt nokkuð mörg svör en það er útfyrir gagnsemi þessarar samlíkingu.
Þegar ýtt er á þennan takka er minnisvæði í reiknivélinni tekið frá og skrifað er í það gildi, sem er svo sótt þegar ANS er notað í útreikningi.
Að sama skapi má láta Python úthluta minnissvæði í tölvunni fyrir þær breytur sem þið viljið geyma.
Munurinn er sá að þið nefnið sjálf hvað minnisvæðið er merkt sem, eruð ekki bundin við að nota ANS og að þið eruð svo gott sem með óteljandi minnissvæði.

Að gefa minnissvæði merkingu og gildi er gert með \textit{gildisveitingu}.
Gildisveiting þýðir að nú er einhver ákveðinn merkimiði kominn með eitthvað til að geyma.
Sjáum einfalt dæmi um þetta.

\begin{lstlisting}[caption=Breytur kynntar]
# Hér er ég að fara að búa til breytu sem heitir val
val = 5

# Þegar ég keyri línuna fyrir ofan segi ég vélinni að hafa aðgengilegt minnisvæði sem ég get notað með því að skrifa orðið val, og settu í það svæði gildið 5.

# Svo ég er að veita breytunni val gildið 5, þess vegar er það kallað gildisveiting.

# Svo get ég notað breytuna mína
# þegar þetta er keyrt fæst svarið 10
val + 5
\end{lstlisting}

Ef þú prófar þig áfram við að búa til breytur gætir þú rekist á svolítið sem hefur ekki gerst áður í vinnubók, að þegar sella inniheldur eingöngu gildisveitingu og er keyrð þá ,,gerist ekkert''.
Þetta finnst mörgum mjög skrýtið því þau vilja fá einhverja útkomu.
En útkoman er sú að þú sagðir vélinni að geyma þetta, þú sagðir henni ekki að gera neitt annað.

Breytur eru skilgreindar vinstra megin við jafnaðarmerki í Python.
Eins og það væri lesið, val fær gildið 5.
Það væri lítið vit í því að hafa það öfugt, 5 er núna jafngilt val.
Það sem við værum þá að segja tölvunni að í hver sinn sem hún vill nota heiltöluna fimm þá á hún að hætta við að nota töluna sjálfa og í staðinn vísa eingöngu í það sem er í minnissvæði merktu val.
Það er alls ekki það sem við viljum.

Nokkrar reglur í nafnavali á breytum, þetta vill vefjast fyrir sumum en lærist fljótlega:
\vspace{0.5cm}
\begin{enumerate}
	\item Kóðalitunin á breytuheitinu má ekki vera annað en venjulegi liturinn fyrir kóða, þannig að ef nafnið fær áherslumerkingu (annan lit) er það ekki löglegt breytuheiti. 
	Áherslulitunin í númeruðu kóðabútunum í þessari bók er marklaus því hún er mjög frumstæð.
	Dæmi um það sem fær áherslulitun eru frátekin lykilorð og tölustafir.
	\item Breytuheitið ætti ekki að innihalda séríslenskan staf (það er löglegt í jupyter vinnubókum en er hrikalega slæmur ávani því það er ekki löglegt allsstaðar).
	\item Breytuheitið má ekki innihalda bil.
\end{enumerate}
\vspace{0.5cm}
Nokkur tilmæli um breytunöfn með tilliti til nafnavenja í Python:
\vspace{0.5cm}
\begin{enumerate}
	\item Breytuheiti byrja á litlum staf.
	\item Ef það þarf að gera löng breytuheiti er venjan að nota snákaframsetningu (e. snake casing) sem felur í sér að gera niðurstrik á milli orða, dæmi \texttt{thetta\_er\_langt\_nafn\_a\_breytu}.
	Annars er til kamelframsetning (e. camel casing) sem felur í sér að annað hvert orð er með stórum staf, dæmi \texttt{thettaErLikaLangtBreytuheiti}.
	Hvort sem þið endið á að nota meira, haldið ykkur bara við annað þeirra.
	\item Breytuheiti eiga að vera lýsandi.
	Ef ég væri að reikna hliðar í þríhyrningi væri gott að eiga breyturnar \texttt{a}, \texttt{b} og \texttt{c}.
	En ef ég væri að búa til reiknirit sem býr til tölvuleikjapersónu af handahófi með því að velja tilviljanakennt nafn, aldur og starf þá væru breytuheitin \texttt{a}, \texttt{b} og \texttt{c} alveg glötuð því þegar ég kæmi aftur að kóðanum mínum myndi ég ekki hafa hugmynd um hvað a, b og c væru. 
	Betra væri að breyturnar hétu \texttt{nafn}, \texttt{aldur} og \texttt{starf}.
\end{enumerate}

\begin{lstlisting}[caption=Dæmi um gildisvetingar\, réttar og rangar]
# Hér er ég að fara að búa til breytu sem heitir val
val = 5

# Hér er ég ekki að búa til breytu sem heitir val heldur er ég að segja að talan fimm er ekki lengur til sem heiltala heldur gæti hún vísað í hvað sem er sem er geymt í minnisvæði merktu val, ólöglegt.
5 = val

# Hér bý ég til breytu sem heitir heiltala sem fær gildið 0
heiltala = 0

# Hér yfirskrifa ég lykilorðið fyrir týpuna heiltala og læt það innihalda 0
int = 0
# þetta er harðbannað og ef þetta gerist er ekki nóg að þurrka þetta út og keyra aftur, nú þarf að endurræsa kjarna vinnubókarinnar.

Gott nafn = 1.0
# Þetta er ekki bara bannað vegna bilsins á milli orðanna ,,Gott'' og ,,nafn'' heldur er það líka ljótt því að það byrjar á stórum staf

3_litlar_mys = 3
# má ekki byrja á tölustaf eða tákni

utreiknud_laun_eftir_skatt = 0.65 * laun
# frábært, lýsandi og gott breytuheiti (hér er þó gert ráð fyrir að vélin þekki breytuna laun)
\end{lstlisting}

Nú þegar við höfum séð hvernig má skilgreina breytu viljum við vita hvernig á að nota þessa breytu.
Ef við snúum okkur aftur að reiknivélasamlíkingunni um ANS takkann þá ætti kóðabútur \ref{lst:notabreytu} að geta sýnt með eðlislægum hætti hvernig breytur nýtast.
Fyrst segi ég vélinni hvað það er sem ANS vísar á, svo segi ég vélinni að mig langar til þess að búa til nýja breytu sem á að byggja á því sem ANS inniheldur.
Í þessum kóðabút er svo haldið áfram með þessa afleiddu breytu og önnur afleidd breyta búin til útfrá henni.
Það sem gerist svo í endann er sambærilegt við það að ýta á ,,='' takkann á reiknivélinni.
Takið eftir að þarna er notuð ný framsetning sem við höfum ekki séð áður, þarna stendur print með svigum fyrir aftan og inni í svigunum er breytan okkar.
Ef þessi kóðabútur er keyrður þá kemur á \textit{staðalúttak} 
\footnote{Þann stað sem texti myndi prentast þegar forritið er notað, hvort sem það er á skjá eða beint á pappír úr prentara eða eitthvað allt annað. Kannski verður úttaki varpað beint inn í heilann á forriturum einhvern tíma?} 
það gildi sem breytan \textbf{x} inniheldur.
Ef þar hefði staðið \texttt{print(halft\_x)} hefðum við fengið svarið sem er geymt í breytunni \texttt{print(halft\_x)}.

\begin{lstlisting}[caption=Að nota breytu, label=lst:notabreytu]
# Hér framkvæmi ég einhvern útreikning sem ég geymi í breytunni ANS
ANS = 5**2 + (4+8.9)**2

# Segjum að þetta hafi verið endapunkturinn í löngu algebrudæmi og nú veit ég hvað y er, og get þá nýtt það til að finna x eins og verða vill svo oft í stærðfræði að x sé týnt. Gefum okkur að x = 3 * y og því fæst
x = 3 * ANS

# Nú ef við viljum reikna eitthvað út með x eigum við það til í minnissvæði merktu x með réttu gildi. Til dæmis með því að búa til breytu fyrir hálft x.
halft_x = x/2


# Nú langar okkur til að vera viss um að við séum við vitrænt svar svo við biðjum tölvuna um að segja okkur hvað er geymt í breytunni x.
print(x)

\end{lstlisting}

Við megum beita print skipuninni óspart og hvetur höfundur til þess að lesandi venji sig á að skoða úttakið sitt í hverju þrepi áður en leitað er hjálpar til annarra.
Print er \textit{fall}, við skoðum föll nánar í kafla \ref{k:föll} en þangað til munum við kynnast nokkrum innbyggðum föllum eins og print().

Núna höfum við séð tvær týpur, heiltölur og fleytitölur.
Breyta getur innihaldið hvernig týpu sem er.
Þá þurfum við að athuga hvað má gera við breyturnar okkar.
Hingað til höfum við eingöngu skoðað reikniaðgerðir sem eru framkvæmdar með kunnuglegum táknum, við höfum ekki verið að beita neinum innbyggðum \textit{aðferðum} á tölurnar okkar.
Við sjáum það gert í kafla \ref{k:strengir} þegar við skoðum hvernig megi vinna með texta.

Að því sögðu þá þurfum við að skoða breytur nokkuð betur áður en við förum að beita þeim á skilvirkan hátt.
Við erum búin að skoða reiknivirkja og gildisveitingu, og nú ætlum við að skoða \textit{reiknivirkjagildisveitingu} þar sem við uppfærum gildi í breytu með því að nota reiknivirkja eins og + eða - með gildisveitingu =.
Þetta sést betur í kóðabút

\begin{lstlisting}[caption="Reiknivirkjagildisveiting", label=lst:reiknivirkjagildisveiting]
# ég ætla að telja nemendur inn í stofuna
# ég byrja með 0 nemendur
nem = 0

# svo sé ég fyrsta nemandann minn og endurskilgreini þar með breytuna nem
nem = nem + 1

# þá uppfæri ég gildið sem nem breytan inniheldur og er hún núna það sem hún var (0) + 1, ef ég keyrði þessa línu aftur væri það orðið að (1) + 1 og svo koll af kolli eftir því sem ég keyri þessa línu oftar

# önnur leið til að skrifa þetta er með reiknivirkjagildisveitingu
nem += 1

# þá uppfæri ég gildið nem um það sem hún var + 1

#segjum að nemendur koma inn í stofuna í pörum þá væri formúlan svona:
nem += 2

# þá uppfærist gildið í nem um +2 í hvert sinn sem þessi kóðalína er keyrð

# Ef ég væri að reikna stofnstærð á bakteríum sem tvöfaldast á hverjum klukkutíma gæti ég gert það svona:
stofn_staerd = 30
stofn_staerd *= 2 # og keyrt svo þessa línu fyrir hvern klukkutíma

# Tökum annað dæmi, byrjum á að skoða hvernig megi geyma útreikninga í mismunandi breytum
thusund = 1000
fimm_hundrud = thusund/2
tvo_hundrud_og_fimmtiu = fimm_hundrud/2

# ef ef mér er sama um breytuna thusund og vil þess í stað bara halda utan um alla upphæðina mína í einni breytu og helminga hana tvisvar þá get ég gert þetta
allt = 1000
allt /= 2
allt /= 2

#nú er allt orðið að 250

# ég má líka nota aðrar breytur í uppfærslunni minni
# hér ætla ég að reikna út hver hækkun launa yrði milli ára ef ég fengi alltaf 2% launahækkun

laun = 100
verdbolga = 0.02
laun *= 1 + verdbolga # þetta má svo keyra endurtekið til að skoða fram í tímann

\end{lstlisting}

Hér sést að það er gagnlegt og fljótlesið þegar það á að uppfæra gildi á breytu að gera það með því að nota reiknivirkjann með gildisveitingunni.
Í kóðabútnum að ofan sést að athugasemdir eru skrifaðar inni í línum líka, það er stundum gagnlegt að skrifa stuttar athugasemdir inni í kóða með þessum hætti en betra er þó að skjala hann skilmerkilega efst við viðeigandi kóðabút.
Allar athugasemdir eru hunsaðar af vélinni og því hefur allt sem er fyrir aftan \# merkið, hver sem það er, engin áhrif á útkomuna.

\section{Villur}\index{Villur}\label{uk:tolur-villur}
Að svo stöddu er gott að hafa í huga hinar ýmsu villur sem geta komið upp.
Þegar við forritum í Python fáum við yfirleitt í hendurnar mjög lýsandi villur sem segja okkur hvað við gerðum vitlaust og hvar.
Áður en lengra er haldið er vitlegt að kynnast þeim algengu villum sem koma upp, hvað þær þýða og hvernig á að taka á þeim.

\begin{itemize}
	\item Nafna villa (e. NameError), þar sem við reyndum að nota breytu sem við höfum enn ekki skilgreint.
	Þessi villa kemur upp þegar við skrifum nafnið á breytunni okkar vitlaust þegar við ætlum að beita henni eða við gleymdum að skilgreina hana áður en við beittum henni.
	Til dæmis myndum við fá nafnavillu við að skilgreina breytuna Valborg en nota svo breytuna valborg, þar sem hástafir og lágstafir skipta máli í Python og eru þessar tvær breytur algerlega óskyldar í minninu.
	Einnig getur þessi villa komið upp þegar við ætlum að nota strengi en gleymum að skilgreina þá sem strengi með gæsalöppum og látum vélina halda að við séum að vinna með breytur.
	Sjáum í kafla \ref{k:strengir} hvernig á að skilgreina strengi.
	
	Þetta lögum við með því að sjá hvaða nafn við ætluðum að nota og pössum okkur að það nafn hefur fengið einhverja skilgreiningu áður en að við reynum að nota það (það er ofar í kóðanum).
	
	\item Týpu villa (e. TypeError), þar sem eitthvað var gert við gögn sem týpan af gögnunum leyfir ekki.
	Til dæmis að beita reiknivirkja á breytu sem styður ekki notkun hans.
	\item Málskipunar villar (e. SyntaxError), þar sem eitthvað var vitlaust skrifað, vitlaust tákn á vitlausum stað.
	Til dæmis var komma notuð í stað punkts eða breytuheiti byrjaði á tákni eða tölu.
	
	Það sem við gerum í því að fá þessa villu er að lesa okkur til um hvað það er sem má gera við týpuna okkar og athuga hvernig við fáum fram þá virkni sem við vildum með því að leysa vandann á annan máta.
	Hér erum við í rauninni að misskilja hvað má og hvað má ekki og við þurfum að átta okkur betur á því hvað er í boði.
	
	\item Inndráttar villa (e. IndendationError), þar sem kóðinn er ekki rétt inndreginn.
	Eins og kom fram í inngangi þá þarf Python kóði að vera vel uppsettur til þess að hann keyri.
	Við sjáum fyrst í kafla \ref{k:segðir} hvernig við beitum inndrætti svo núna ef við fáum þessa villu þá er það vegna þess að við erum með óþarfa bil fyrir framan kóðann okkar. 
	
	Það sem við gerum er að skilja hvað á að vera í hvaða inndrætti og laga kóðann þannig að hann sé vel uppsettur.
	
	\item Vísis villa (e. IndexError), þar sem við reyndum að ná í sætisvísi sem er ekki til í gögnunum okkar.
	Í kafla \ref{k:strengir} sjáum við týpu sem hefur sætisvísa.
	
	Það sem við gerum í þessari villu er að átta okkur á því hversu margir sætisvísar eru til staðar og hvaða vísa við vildum fá, laga svo kóðann í samræmi við það sem við ætluðum að sækja.
	
	\item Gildis villa (e. ValueError), þar sem við reyndum að ná í eða nota gildi úr gögnum sem er ekki til.
	Gögnin eru til en gildið er ekki til staðar í þeim.
	
	Athugum hér að við reyndum að sækja eitthvað, næstum því eins og upp úr poka, og í stað þess að segja áttu þetta til? Þá sögðum við við vélina ,,þú átt þetta til, láttu mig hafa það!''. Svo það sem við getum gert til að laga það er að segja ,,áttu þetta til?'' eða ,,reyndu að láta mig hafa það ef þú átt það''. 
	Þessi lausn er ansi óljós að svo stödu þar sem við höfum ekki farið yfir hvernig á að ná í gildi upp úr gögnunum okkar en vonandi verður þetta ljósara þegar þar að kemur.
	
	\item Eiginda villa (e. AttributeError), þar sem við reynum að ná í eða beita eigindum sem gögnin okkar búa ekki yfir.
	Við getum séð hvaða eigindum gögnin okkar búa yfir með því að lesa skjölun um þau eða kíkja snöggt undir húddið með því að gera punkt fyrir aftan gögnin og ýta á tab takkann \todo{vísa hérna aftur í tab takkann} og þá sjáum við þau eigindi sem týpan býr yfir.
	Ef við skoðum þau eigindi sem eru til fyrir heiltölur og fleytitölur sjáum við að þau eru ekki alveg eins og ef við reynum að nota þau sem eru til fyrir annað en ekki hitt á víxl þá fáum við eigindavillu.
	
	Það sem við gerum í þessu er svipað týpu villunni, það er að lesa okkur til um hvaða eigindi gögnin okkar búa yfir og hvað við getum náð í og notað.
	Eigindi eru mismunandi eftir týpum og er mikilvægt að átta sig á því hvaða eigindi eru í boði hverju sinni svo að við veljum rétta týpu fyrir gögnin okkar.
\end{itemize}

\section{Verkefni}
\begin{exercise}\label{tol1}
		Búðu til breytu sem inniheldur heiltölu
	\end{exercise}
\setboolean{firstanswerofthechapter}{true}
	\begin{Answer}[ref={tol1}]
		Nú þurfum við að athuga þrennt nafn á breytuna, hvernig á að veita henni eitthvað gildi og sjá til þess að það gildi sé heiltala.
		Til þess að búa til breytu þarf að finna nafn á hana, skrifa það fremst í línuna (ss. vinstra megin við jafnaðarmerkið) og nota svo jafnaðarmerkið til þess að veita breytunni gildi og því næst velja einhverja heila tölu.
		Það er tölu sem er annað hvort jákvæð eða neikvæð og er ekki með punkt og aukastafi.
\begin{lstlisting}
# til dæmis
eitthvad_loglegt_breytuheiti_fyrir_heiltolu = 5
neikvaed_tala = -1234567890\end{lstlisting}
\end{Answer}
\setboolean{firstanswerofthechapter}{false}

\begin{exercise}\label{tol2}
	Búðu til breytu sem inniheldur fleytitölu
\end{exercise}
\begin{Answer}[ref={tol2}]
	Hér þurfum við að hafa það sama í huga eins og í æfingu \ref{tol1} nema að talan sem við veljum okkur verður að vera með punkti og einum eða fleiri aukastöfum.
\begin{lstlisting}
# til dæmis
eitthvad_loglegt_breytuheiti_fyrir_fleytitolu = 5.0
neikvaed_tala = -12.9999\end{lstlisting}
\end{Answer}

\begin{exercise}\label{tol2-5}
	Búðu til breytu sem inniheldur samlagningu breytanna úr verkefnum \ref{tol1} og \ref{tol2}. Af hvaða týpu er þessi þriðja breyta?
\end{exercise}
\begin{Answer}[ref={tol2-5}]
	Nú þurfum við að athuga að nota nákvæmlega sömu nöfn og við notuðum í fyrri verkefnum.
	Vegna þess að hér erum við að vísa í breytur sem búið var að skilgreina, við fáum nafnavillu ef við skrifum nöfn hægra megin við jafnaðarmerkið sem tölvan hefur ekki aðgang að í minni.
\begin{lstlisting}
# til dæmis
thridja_breytan = eitthvad_loglegt_breytuheiti_fyrir_heiltolu + eitthvad_loglegt_breytuheiti_fyrir_fleytitolu
print(thridja_breytan)\end{lstlisting}
	Svo er einnig spurt af hvaða týpu er þessi þriðja breyta?
	Þar sem önnur breytan er fleytitala verður útkoman fleytitala, það má sjá með því að gera eftirfarandi:
	
\begin{lstlisting}
type(thridja_breytan)\end{lstlisting}
\end{Answer}


\begin{exercise}\label{tol3}
	Gefum okkur að við eigum eftirfarandi kóða, hvernig fáum við breytuna helmingur til að innihalda helminginn af því sem er í breytunni allt.
	Nú megum við ekki bara skrifa inn töluna 1000 í útreikninginn, við þurfum að beita breytunni \textit{allt} því að hún gæti breyst og við viljum ekki að kóðinn okkar gefi vitlaust svar ef \textit{allt} breytist.
	Einnig viljum við búa til þriðju breytuna sem á að vera helmingurinn af helmingnum, hvernig fáum við helminginn af breytunni \textit{helmingur}.
\begin{lstlisting}
allt = 1000
helmingur =\end{lstlisting}
\end{exercise}
\begin{Answer}[ref={tol3}]
	Við beitum deilingu á breytuna allt, og svo gerum við það aftur fyrir nýju breytuna.
\begin{lstlisting}
allt = 1000
helmingur = allt/2
helmingurinn_af_helmingnum = helmingur/2

# eða
allt = 1000
helmingur = allt/2
helmingur_af_helmingnum = allt/2/2  # eða allt/4 því helmginur af helmingnum er fjórðungur\end{lstlisting}
\end{Answer}


\begin{exercise}\label{tol4}
	Nú langar okkur að reikna helminginn af breytunni \textit{allt} alveg eins og í verkefni \ref{tol3} nema við viljum alls ekki geyma niðurstöðurnar í breytu heldur viljum við hafa áhrif á breytuna \textit{allt}.
	Hvernig getum við prentað út fyrst helminginn af \textit{allt} og svo helminginn af því?
\end{exercise}
\begin{Answer}[ref={tol4}]
	Við beitum deilingu á breytuna allt, og svo gerum við það aftur fyrir nýju breytuna.
\begin{lstlisting}
allt = 1000
print(allt)

allt /= 2
print(allt)

allt /= 2
print(allt)\end{lstlisting}
\end{Answer}

\begin{exercise}\label{tol5}
	Hvernig fyllirðu inn í eftirfarandi kóða þannig að hann reikni út langhlið í rétthyrndum þríhyrningi?
	Þar sem a og b eru skammhliðar þríhyrningsins og þær þarf að skilgreina sem einhverjar stærðir áður en hægt er að reikna c.
\begin{lstlisting}
??
??

c_i_odru_veldi = a**2 + b**2
c = c_i_odru_veldi**0.5
print(c)\end{lstlisting}
\end{exercise}
\begin{Answer}[ref={tol5}]
	Hér þurfum við að átta okkur á því að vélin býst við að nota breyturnar a og b en við eigum eftir að skilgreina þær sem eitthvað.
	Þekktur rétthyrndur þríhyrningur er til dæmis með hliðarnar 3, 4 og 5.
\begin{lstlisting}
a = 3
b = 4

c_i_odru_veldi = a**2 + b**2
c = c_i_odru_veldi**0.5
print(c)\end{lstlisting}
	Þá er c 5 í þessu tilfelli en hér má setja hvaða tölur inn sem er til að komast að því hvað c þyrfti að vera löng hlið til þess að þríhyrningurinn sé rétthyrndur.
\end{Answer}

\begin{exercise}\label{tol-villa}
	Skrifaðu kóða sem veldur eftirfarandi villum:
	\begin{tasks}(2)
		\task\label{tol-villa1} Nafnavilla.  
		\task\label{tol-villa2} Málskipunarvilla.
		\task\label{tol-villa3} Inndráttarvilla.
	\end{tasks}
\end{exercise}
\begin{Answer}[ref={tol-villa}]
	Þetta eru tillögur, en alls ekki eina leiðin til að fá þessar villur.
\begin{lstlisting}
# nafnavilla
breyta = breyta8 # breyta8 hefur ekki verið skilgreind

# Málskipunarvilla
3breyta = 3 # tölustafur má ekki vera fremst í breytuheiti

# Inndráttarvilla
breyta3 = 3 # of mikið bil fremst í línunni\end{lstlisting}
\end{Answer}

\chapter{Strengir}\index{Strengir}\label{k:strengir}
Til þess að geta sýnt og notað texta þarf gagnatýpu til að halda utan um hann. Í flestum forritunarmálum, og Python er ekki undantekning, eru gögn af þeirri týpu kölluð \textbf{strengir}.
Lykilorð fyrir þessa týpu er \textbf{str}.

\section{Strengir skilgreindir}\index{Strengir skilgreindir}
Til þess að afmarka texta og segja vélinni að fara með hann sem af týpunni strengur þarf að nota tákn.
Við þurftum ekki að gera það þegar við skrifuðum tölurnar en nú, og seinna, munum við þurfa sér tákn til þess að segja vélinni gögn af hvaða týpu hún er að vinna með.

Táknin sem skilgreina strengi eru gæsalappir.
Einfaldar eða tvöfaldar.

\begin{lstlisting}[caption="Strengir skilgreindir]
# Fyrsti strengurinn okkar
"halló"

#strengur geymdur í breytu
textinn_minn = "halló ég má skrifa mörg orð inn í þessar gæsalappir"

einfaldar_gaesalappir = 'ég má líka skrifa innan einfaldra gæsalappa'

thetta_virkar_ekki = 'gæsalappirnar þurfa að passa saman" 

# og ef ég vil skrifa mjög langan texta nota ég þrjár gæsalappir
langi_textinn_minn = ''' ég má skrifa eins langa setningu hér og ég vil því að þetta verður alltaf álitið sem ein lína, hins vegar ef ég nota öðruvísi gæsalappir og langar að gera kóðan læsilegan er hægt að brjóta hann upp án þess að nota þessar þreföldugæsalappir, við sjáum það eftir smá'''
\end{lstlisting}

Í sumum forritunarmálum er munur á því að nota einfaldar og gæsalappir, þar sem einfaldar eru notaðar fyrir staka stafi (sér gagnatýpa) og tvöfaldar fyrir strengi.
En það er enginn raun munur á því hvernig Python meðhöndlar þær.

\section{Strengir og reikniaðgerðir}\index{Strengir og reikniaðgerðir}
Við erum búin að sjá að það megi leggja tölur saman og margfalda þær.
Nú ætlum við að skoða hvaða reikniaðgerðir er hægt að framkvæma með strengi og hvaða áhrif það hefur.

\todo{vísa í vinnubók inni í þessum kóðabút}

\begin{lstlisting}[caption="Strengir og reikniaðgerðir"]

# Reikniaðgerðirnar sem við þekkjum eru +, -, *, /, //, **, og %
# Lesandinn er hvattur til þess að gera prófanir á þessu í vinnubók upp á eigin spýtur
# Með því að skilgreina streng og reyna að nota reikniaðgerð á hann með tölum eða öðrum strengjum


# Gerum ráð fyrir að þessar prófanir hafi átt sér stað og niðurstaðan sé sú að þær aðgerðir sem hægt er að framkvæma eru + og *
# En hvað gerist þegar við notum þær?

\end{lstlisting}


Þegar við notum + til að setja saman strengi þá erum við að beita \textit{samskeytingu} (e. concatenation).
Samskeyting þýðir að einum streng er bætt við fyrir aftan annan streng.
Það skiptir máli hvor er fyrir framan: "halló" + "bless" verður að "hallóbless" en "bless" + "halló" verður að "blesshalló".

\begin{lstlisting}[caption="Samskeyting strengja"]

strengur_a = "þetta er a strengurinn minn"
strengur_b = " og þetta er b strengurinn minn"

# Nú get ég sameinað þessa strengi með því að setja annan þeirra fyrir aftan hinn
sameinadir_a_og_b = a + b

# ef við prentum út strenginn fáum við 
"þetta er a strengurinn minn og þetta er b strengurinn minn"
#takið eftir að það er bil á milli strengjanna, það er eingungis vegna þess að b strengurinn er skilgreindur þannig að fyrst kemur bil fremst í strengnum

# röðin skiptir máli þegar strengir eru sameinaðir svona
sameinadir_b_og_a = b + a

# þetta útprentað skilar okkur:
" og þetta er b strengurinn minnþetta er a strengurinn minn"

# Takið eftir að þarna er ekkert bil á milli strengjanna.

# Á þessu er hægt að svindla:
fyrsta_nafn = "Valborg"
seinna_nafn = "Sturludóttir"
fullt_nafn = fyrsta_nafn + " " + seinna_nafn

# Þarna sameinaði ég þrjá strengi þar sem ég vissi að hvorugur strengjanna minna innihéldi bil ákvað ég að setja það á milli með auka samskeytingu.
\end{lstlisting}

Þegar við notum * til þess að margfalda streng erum við að \textit{lengja} (e. multiply) hann.
Strengjalenging virkar þannig að þú tilgreinir hversu oft, í heilum tölum, þú vilt að strengurinn sé endurtekinn.

\begin{lstlisting}[caption="Strengjalenging"]
eitt_ord = "kex"
eitt_ord*3

#skilar okkur
"kexkexkex"
\end{lstlisting}

\section{Vísar í streng}\index{Vísar í streng}
Strengir eru af einhverri lengd, við getum séð hversu mörg stafabil eru í streng með því að telja þau sjálf eða láta tölvuna segja okkur það með innbyggða fallinu len() (fyrir length).

\begin{lstlisting}[caption="Stafabilafjöldi", label=lst:stafabil]
strengur1 = "kex"
# þessi strengur er augljóslega þrjú stafabil
len(strengur1) 
# skilar okkur 3

strengur2 = "kex með smjöri, osti og sultu"
len(strengur2)
# skilar okkur 29, takið eftir að tóm bil og greinamerki teljast með
\end{lstlisting}

Nú þegar við vitum hversu mörg stafabil eru í strengnum getum við notað þau.
Við getum sagt við vélina mig langar til að fá vísi (e. index) (einnig kallað sætisnúmer, sæti og stæði) númer 1 og séð hvaða tákn er í þeim vísi.
Til að ná í eitthvað upp úr streng þurfum við að nota hornklofa (e. square brackets), tákn sem eru eins og kassalaga svigar [ og ].
Við notum þessi tákn í Python til að ná í gögn upp úr einhverri gagnagrind, sjáum nánari útskýringu á því fyrirbæri í kaflanum \ref{k:listar}.
Nú lítum við svo á að strengir séu til þess að geyma fyrir okkur tákn í ákveðinni röð og við getum nálgast þessi tákn með því að nota hornklofa.
Inn í hornklofann ætlum við að láta þann vísi (eða það sætisnúmer) sem við viljum vinna með.

\begin{lstlisting}[caption="Vísir 1", label=lst:stafbil2]
strengur = "kex"
# þessi strengur er augljóslega þrjú stafabil

print(strengur[1])

# þetta skilar okkur stafnum 'e'
\end{lstlisting}

Eins og sést í kóðabút \ref{lst:stafabil2} þá vísar vísir númer 1 ekki á fremsta stafinn sem er í þessu tilfelli k heldur vísar hann á stafinn e.
Það er vegna þess að í Python og flestum öðrum forritunarmálum (ekki öllum) er byrjað að telja í núll.
Þannig að fremsti vísirinn í streng (og öðrum gagnagrindum) er núllti vísirinn.
Hver er þá síðasti vísirinn?
Nú höfum við komist að þeirri niðurstöðu (í kóðabút \ref{lst:stafabil}) að strengurinn "kex" hefur þrjú stafabil, að það séu þrír sætisvísar í strengnum, að k sé í vísi 0, e sé í vísi 1 og þá hlýtur x að vera í vísi 2.
Síðasti vísirinn í streng er því einum lægri heldur en lengdin á strengnum.
Þannig að strengur af lengdinni fimm, eins og strengurinn "texti", hefur fimm stafabil sem eru í vísum númer 0,1,2,3 og 4.

\subsection{Óbreytanleiki}
Nú höfum við séð að það er hægt að sækja stafabil í streng, eins og tildæmi núllt táknið í strengnum.
Þá er mikilvægt að hafa í huga að í Python er ekki leyfilegt að endurskilgreina hluta úr streng.
Byrjum á því að skoða hvað endurskilgreining þýðir.
Ef við búum til breytu eins og í kóðabút \ref{lst:reiknivirkjagildisveiting} og notum nafnið á henni aftur til að segja vélinni að endurnýta minnissvæði með ákveðnu nafni.
Þá erum við búin að endurskilgreina breytuna okkar, hún var eitthvað áður en nú er hún eitthvað annað.

Þar sem strengir eru með ákveðin númeraða vísa sem benda á ákveðin tákn gætum við þá ekki bara sagt við vélina ,,mig langar að breyta bara tákni númer 0''?
Það er ekki í boði því að í Python eru strengir óbreytanlegir (e. immutable) og því er bara hægt að vinna með því eins og þeir eru eða endurskilgreina þá alveg.

\subsection{Neikvæðir vísar}
Það má einnig telja frá hægri til vinstri.
Ef okkur langar að vinna með öftustu stökin í streng þurfum við ekki að vita hvað strengurinn er langur, við getum talið frá hægri endanum og unnið með neikvæða vísa.
Í því tilfelli byrjum við ekki að telja í 0, því að þá myndi verða til tvíræðni (e. ambiguity).
Tölvan myndi ekki vita hvorn 0 vísinn við værum að biðja um þegar við segðum strengur[0], hvort við værum að tala um núll frá vinstri eða hægri.
Þess vegna byrjum við að telja frá hægri í -1, og höldum þannig áfram þar til við erum komin niður í -n þar sem n er lengdin á strengnum.
Svo strengurinn "kex" er með vísana 0,1 og 2 en einnig vísana -3, -2 og -1 bæði í þessari röð, svo vísir -1 er alltaf síðasta táknið í streng.

\subsection{Hlutstrengir}
Nú vitum við hvernig á að sækja eitt stakt tákn upp úr streng.
En hvernig náum við í einhvern hluta úr honum?
Það er einnig gert með hornklofunum og við notum þá með ákveðnum hætti, við fáum að setja inn fleiri upplýsingar heldur en bara hvaða staka vísi við viljum.
Við notum vísana svona: [vísir sem á að byrja fyrir framan : vísir sem á að hætta fyrir framan: stærð á skrefi sem á að taka við lesturinn].
Þarna eru semsagt tveir tvípunktar sem er verið að nota og tölurnar sem koma á milli þeirra er afmörkunin á því hvað við viljum lesa upp úr tilteknum streng.
Nú er vert að nefna að þegar við notum þessa málskipan eru ákveðin gildi sjálfgefin, það er við þurfum ekki að taka þau fram.

Sjálfgefin gildi við að ná í hlutstreng:
\begin{enumerate}
	\item vísirinn sem við byrjum fyrir framan er fremsti stafurinn, eða vísir 0
	\item vísirinn sem við hættum fyrir framan er fyrir aftan aftasta stafinn, ekki vísi -1 heldur fyrir aftan hann, enginn vísir
	\item skrefastærðin er sjálfgefin 1, það er að við skoðum hvert einasta tákn og hoppum ekki yfir neitt
\end{enumerate}

\begin{lstlisting}[caption="Hlutstrengir", label=lst:hlutstrengir]
strengur = "kex með smjöri, osti og sultu"

# mig langar að sækja allan strenginn
sami_strengur = strengur[::] # hér nota ég þann kost í Python að setja inn fyrir mig sjálfgefin gildi

#mig langar að sækja það sem kemur fyrir aftan táknið x í kex
aftan_x = strengur[3::]

#mig langar ekki í síðasta stafinn
nema_sidasti = strengur[:-1:]

#mig langar í allan strenginn nema annan hvern staf
annar_hver = strengur[::2]

# mig langar að sækja orðið kex og mig langar einnig að sækja orðið sultu
kex = strengur[0:3]
sultu = strengur[-5:]

#einnig er hægt að lesa afturábak með því að taka neikvæða skrefastærð
ofugur = strengur[::-1]
\end{lstlisting}


\section{Aðferðir á strengi}\index{Aðferðir á strengi}
Áður en aðferðir á strengi eru kynntar þarf að útskýra stuttlega hvað aðferðir eru.
Við höfum séð print() fallið notað t.d. í kóðabút \ref{lst:notabreytu} og í kóðabút \ref{lst:stafabil}, það er innbyggt fall í Python sem prentar það sem beðið er um á staðalúttak. \todo{er ég búin að útskýra staðalúttak?}
Það fall virkar eitt og sér og þarf bara að skrifa nafnið á því, gera sviga og setja inn í svigana það sem við viljum láta fallið fá \todo{vísa í falla kaflann}.
Aðferðir eru sérhæfð föll sem virka á ákveðnar gagnatýpur.
Þannig að allar aðferðir eru föll, ekki öll föll eru aðferðir.
Á ensku eru aðferðir kallaðar \textit{methods} og föll \textit{functions}.
Aðferðir eru í raun ,,hengdar aftan á" þá týpu sem þær eiga að verka á, það er gert með því að skrifa nafnið á breytunni sem inniheldur gögnin sem við viljum framkvæma aðferðina á, gera svo punkt, skrifa nafnið á aðferðinni og setja sviga, inn í svigana fara öll þau viðföng sem aðferðin tekur við.
\todo{vísa í eitthvað sem útskýrir viðföng}

\begin{lstlisting}[caption="Aðferðir kynntar", label=lst:adferdir]
strengur = "kex með smjöri, osti og sultu"
\end{lstlisting}

Annað sem þarf að hafa í huga áður en við vinnum með aðferðir á strengi er að strengir eru óbreytanlegir (e. immutable) sem þýðir að aðferðir sem eru notaðar á þá skila öðrum strengjum í stað þess að breyta strengnum sem við keyrðum aðferðina á.
Með það í huga skulum við skoða eftirfarandi lista af aðferðum sem áhugavert er að taka fyrir.
\vspace{10px}

Hér koma fyrir nokkrar aðferðir, gerum ráð fyrir að þær séu að verka á breytuna strengur sem inniheldur táknin "valborg Sturludóttir".
\vspace{5px}
\begin{enumerate}
	\item strengur.capitalize() skilar strengnum "Valborg sturludóttir" þar sem fremsti táknið er nú hástafur.
	\item strengur.upper() skilar strengnum "VALBORG STURLUDÓTTIR" þar sem allir stafir eru nú háfstafir.
	\item strengur.lower() skilar strengnum "valborg sturludóttir" þar sem allir stafir eru nú lágstafir.
	\item strengur.switchcase() skilar strengnum "VALBORG sTURLUDÓTTIR" þar sem búið er að skipta út lágstöfum fyrir hástafi og öfugt.
	\item strengur.index('v') skilar tölunni 0 þar sem fyrsta 'v' táknið kemur fyrir í vísi 0
	\item strengur.index('x') skilar villu þar sem táknið 'x' hefur engann vísi í strengnum
	\item strengur.find('v') skilar tölunni 0 þar sem fyrsta 'v' táknið kemur fyrir í vísi 0
	\item strengur.find('x') skilar tölvunni -1 þar sem 'x' finnst ekki í stregnum. 
\end{enumerate}
\vspace{10px}

Takið eftir því að þarna er orðið lykilorðið ,,skilar" það er að við fáum í hendurnar eitthvað til að vinna með sem við getum t.d. vistað í breytu, við skoðum þetta nánar þegar við gerum okkar eigin föll í kafla \ref{k:föll}.
Það er þörf á því að vinna með aðferðir á strengi með þessum hætti því að við munum að strengir eru óbreytanlegir.
Þannig að ef við viljum vinna með einhverja útkomu byggða á streng þá þurfum við að fá útkomuna í hendurnar, því strengurinn sem aðferðinni var beitt á breytist ekki neitt við að kalla í aðferðina.
Í upptalningunni hér að ofan getum við keyrt allar þessar línur í röð eins og kóða og búist við að fá þessi svör því að breytan strengur verður aldrei fyrir neinum breytingum.



% ---------------------------------------------- 

% SNIÐMÁT

% ----------------------------------------------

\subsection{Descriptions and Definitions}\index{Lists!Descriptions and Definitions}

\begin{description}
\item[Name] Description
\item[Word] Definition
\item[Comment] Elaboration
\end{description}

%----------------------------------------------------------------------------------------
%	CHAPTER 2
%----------------------------------------------------------------------------------------

\chapter{In-text Elements}

\section{Theorems}\index{Theorems}

This is an example of theorems.

\subsection{Several equations}\index{Theorems!Several Equations}
This is a theorem consisting of several equations.

\begin{theorem}[Name of the theorem]
In $E=\mathbb{R}^n$ all norms are equivalent. It has the properties:
\begin{align}
& \big| ||\mathbf{x}|| - ||\mathbf{y}|| \big|\leq || \mathbf{x}- \mathbf{y}||\\
&  ||\sum_{i=1}^n\mathbf{x}_i||\leq \sum_{i=1}^n||\mathbf{x}_i||\quad\text{where $n$ is a finite integer}
\end{align}
\end{theorem}

\subsection{Single Line}\index{Theorems!Single Line}
This is a theorem consisting of just one line.

\begin{theorem}
A set $\mathcal{D}(G)$ in dense in $L^2(G)$, $|\cdot|_0$. 
\end{theorem}

%------------------------------------------------

\section{Definitions}\index{Definitions}

This is an example of a definition. A definition could be mathematical or it could define a concept.

\begin{definition}[Definition name]
Given a vector space $E$, a norm on $E$ is an application, denoted $||\cdot||$, $E$ in $\mathbb{R}^+=[0,+\infty[$ such that:
\begin{align}
& ||\mathbf{x}||=0\ \Rightarrow\ \mathbf{x}=\mathbf{0}\\
& ||\lambda \mathbf{x}||=|\lambda|\cdot ||\mathbf{x}||\\
& ||\mathbf{x}+\mathbf{y}||\leq ||\mathbf{x}||+||\mathbf{y}||
\end{align}
\end{definition}

%------------------------------------------------

\section{Notations}\index{Notations}

\begin{notation}
Given an open subset $G$ of $\mathbb{R}^n$, the set of functions $\varphi$ are:
\begin{enumerate}
\item Bounded support $G$;
\item Infinitely differentiable;
\end{enumerate}
a vector space is denoted by $\mathcal{D}(G)$. 
\end{notation}

%------------------------------------------------

\section{Remarks}\index{Remarks}

This is an example of a remark.

\begin{remark}
The concepts presented here are now in conventional employment in mathematics. Vector spaces are taken over the field $\mathbb{K}=\mathbb{R}$, however, established properties are easily extended to $\mathbb{K}=\mathbb{C}$.
\end{remark}

%------------------------------------------------

\section{Corollaries}\index{Corollaries}

This is an example of a corollary.

\begin{corollary}[Corollary name]
The concepts presented here are now in conventional employment in mathematics. Vector spaces are taken over the field $\mathbb{K}=\mathbb{R}$, however, established properties are easily extended to $\mathbb{K}=\mathbb{C}$.
\end{corollary}

%------------------------------------------------

\section{Propositions}\index{Propositions}

This is an example of propositions.

\subsection{Several equations}\index{Propositions!Several Equations}

\begin{proposition}[Proposition name]
It has the properties:
\begin{align}
& \big| ||\mathbf{x}|| - ||\mathbf{y}|| \big|\leq || \mathbf{x}- \mathbf{y}||\\
&  ||\sum_{i=1}^n\mathbf{x}_i||\leq \sum_{i=1}^n||\mathbf{x}_i||\quad\text{where $n$ is a finite integer}
\end{align}
\end{proposition}

\subsection{Single Line}\index{Propositions!Single Line}

\begin{proposition} 
Let $f,g\in L^2(G)$; if $\forall \varphi\in\mathcal{D}(G)$, $(f,\varphi)_0=(g,\varphi)_0$ then $f = g$. 
\end{proposition}

%------------------------------------------------

\section{Examples}\index{Examples}

This is an example of examples.

\subsection{Equation and Text}\index{Examples!Equation and Text}

\begin{example}
Let $G=\{x\in\mathbb{R}^2:|x|<3\}$ and denoted by: $x^0=(1,1)$; consider the function:
\begin{equation}
f(x)=\left\{\begin{aligned} & \mathrm{e}^{|x|} & & \text{si $|x-x^0|\leq 1/2$}\\
& 0 & & \text{si $|x-x^0|> 1/2$}\end{aligned}\right.
\end{equation}
The function $f$ has bounded support, we can take $A=\{x\in\mathbb{R}^2:|x-x^0|\leq 1/2+\epsilon\}$ for all $\epsilon\in\intoo{0}{5/2-\sqrt{2}}$.
\end{example}

\subsection{Paragraph of Text}\index{Examples!Paragraph of Text}

\begin{example}[Example name]
\lipsum[2]
\end{example}

%------------------------------------------------

\section{Exercises}\index{Exercises}

This is an example of an exercise.

\begin{exercise}
This is a good place to ask a question to test learning progress or further cement ideas into students' minds.
\end{exercise}

%------------------------------------------------

\section{Problems}\index{Problems}

\begin{problem}
What is the average airspeed velocity of an unladen swallow?
\end{problem}

%------------------------------------------------

\section{Vocabulary}\index{Vocabulary}

Define a word to improve a students' vocabulary.

\begin{vocabulary}[Word]
Definition of word.
\end{vocabulary}

%----------------------------------------------------------------------------------------
%	PART
%----------------------------------------------------------------------------------------

\part{Part Two}

%----------------------------------------------------------------------------------------
%	CHAPTER 3
%----------------------------------------------------------------------------------------

\chapterimage{chapter_head_1.pdf} % Chapter heading image

\chapter{Presenting Information}

\section{Table}\index{Table}

\begin{table}[h]
\centering
\begin{tabular}{l l l}
\toprule
\textbf{Treatments} & \textbf{Response 1} & \textbf{Response 2}\\
\midrule
Treatment 1 & 0.0003262 & 0.562 \\
Treatment 2 & 0.0015681 & 0.910 \\
Treatment 3 & 0.0009271 & 0.296 \\
\bottomrule
\end{tabular}
\caption{Table caption}
\label{tab:example} % Unique label used for referencing the table in-text
%\addcontentsline{toc}{table}{Table \ref{tab:example}} % Uncomment to add the table to the table of contents
\end{table}

Referencing Table \ref{tab:example} in-text automatically.

%------------------------------------------------

\section{Figure}\index{Figure}

\begin{figure}[h]
\centering\includegraphics[scale=0.5]{placeholder.jpg}
\caption{Figure caption}
\label{fig:placeholder} % Unique label used for referencing the figure in-text
%\addcontentsline{toc}{figure}{Figure \ref{fig:placeholder}} % Uncomment to add the figure to the table of contents
\end{figure}

Referencing Figure \ref{fig:placeholder} in-text automatically.

%----------------------------------------------------------------------------------------
%	BIBLIOGRAPHY
%----------------------------------------------------------------------------------------

\chapter*{Bibliography}
\addcontentsline{toc}{chapter}{\textcolor{ocre}{Bibliography}} % Add a Bibliography heading to the table of contents

%------------------------------------------------

\section*{Articles}
\addcontentsline{toc}{section}{Articles}
\printbibliography[heading=bibempty,type=article]

%------------------------------------------------

\section*{Books}
\addcontentsline{toc}{section}{Books}
\printbibliography[heading=bibempty,type=book]

%----------------------------------------------------------------------------------------
%	INDEX
%----------------------------------------------------------------------------------------

\cleardoublepage % Make sure the index starts on an odd (right side) page
\phantomsection
\setlength{\columnsep}{0.75cm} % Space between the 2 columns of the index
\addcontentsline{toc}{chapter}{\textcolor{ocre}{Index}} % Add an Index heading to the table of contents
\printindex % Output the index

%----------------------------------------------------------------------------------------

\end{document}
